\begin{textsample}{POS Dim 3 – 2014\_c6 – Score 5.00 – 2014\_c6/enade\_2014\_pedag\_076bc59f40.txt}  \label{ex:f3_pos_037}
CAPÍTULO 3 ANÁLISE TÉCNICA DA PROVA Este capítulo tem por objetivo apresentar o desempenho dos estudantes concluintes de Pedagogia ( Licenciatura ) no ENADE/2014.
Para isso, foram calculadas as estatísticas básicas da prova em seu todo ( seção 3.1.1 ), bem como as estatísticas relacionadas aos Componentes de Formação Geral ( seção 3.1.2 ) e de Conhecimento Específico da Área ( seção 3.1.3 ).
Dadas as suas características, foram analisadas, em separado, as questões objetivas ( seção 3.2 ) e as questões discursivas ( seção 3.3 ).
Para as questões objetivas, foram disponibilizados os índices de facilidade e de discriminação bisserial, também em separado, para os Componentes de Formação Geral ( seção 3.2.1 ) e de Conhecimento Específico ( seção 3.2.2 ).
De cada componente, uma das questões foi escolhida para exemplificar a análise gráfica, relacionando as alternativas escolhidas pelos estudantes ( inclusive o gabarito ) com o número de acertos no componente.
O Anexo I apresenta a íntegra da análise gráfica para todas as questões objetivas.
Para cada uma das questões discursivas, os conteúdos dos tipos mais comuns de respostas dos estudantes são apresentados e comparados com o padrão de respostas esperado ( ver Anexo VIII com o padrão de respostas ).
Tomando como base as duas questões discursivas do Componente de Formação Geral, a seção 3.3.1.6 apresenta comentários sobre a correção das respostas com respeito à Língua Portuguesa.
Na Tabela 3.1 são apresentados o tamanho da população inscrita e de presentes, e as seguintes estatísticas das notas18 : média do desempenho na prova, erro padrão da média, desvio padrão, nota mínima, mediana e nota máxima.
Tais estatísticas contemplam o total de estudantes concluintes da Área de Pedagogia ( Licenciatura ) inscritos e presentes à prova do ENADE/2014, tendo em vista as agregações por Grandes Regiões e o país como um todo.
As estatísticas das demais tabelas ímpares serão baseadas na mesma população da Tabela 3.1 e, portanto, do mesmo tamanho.
As três primeiras tabelas pares são desagregadas por Categoria Administrativa e Organização Acadêmica da Instituição de Ensino Superior ( IES ).
As estatísticas dessas tabelas pares serão baseadas na mesma população da Tabela 3.2.
Em relação aos gráficos de distribuição de notas, o intervalo considerado foi de 10 unidades, aberto à esquerda e fechado à direita, com exceção do primeiro intervalo, [ 0 ; 10 ], fechado em ambos os extremos.
Para os gráficos de distribuição das 18 Uma definição dessas estatísticas pode ser encontrada no Glossário. notas das questões discursivas, foram consideradas mais duas categorias : questão em branco19 e nota zero. 3.1 ESTATÍSTICAS BÁSICAS DA PROVA Esta seção apresenta estatísticas selecionadas e histogramas da nota geral ( 3.1.1 ) e de cada componente : Formação Geral ( 3.1.2 ) e Conhecimento Específico ( 3.1.3 ).
São também apresentadas estatísticas selecionadas de subpopulações caracterizadas por Grande Região, Categoria Administrativa e Organização Acadêmica. 3.1.1 Estatísticas Básicas Gerais A Tabela 3.1 apresenta as estatísticas básicas da prova, por grande Região, dos estudantes concluintes de Pedagogia ( Licenciatura ).
A população total de inscritos foi de 131.146.
Destes, 111.863 estiveram presentes, sendo 14,7 \% o índice de não comparecimento.
A Região de maior abstenção foi a Norte ( 22,3 \% ) e a de menor abstenção foi a Nordeste ( 12,6 \% ).
A média das notas da prova como um todo ( nas seções seguintes serão analisados os Componentes de Formação Geral e de Conhecimento Específico ) foi 46,7, sendo que os estudantes da região Centro-Oeste obtiveram a média mais baixa ( 44,2 ), e os da região Sudeste obtiveram a média mais alta ( 48,2 ).
As demais médias foram : 44,4 na região Norte ; 45,6 na região Nordeste ; e 46,6 na região Sul.
O desvio padrão para o Brasil como um todo foi 15,0, sendo o maior desvio padrão encontrado na região Sudeste ( 15,2 ) e o menor, na região Norte ( 14,6 ), indicando uma dispersão um pouco menor das notas desta última região.
A região que obteve a maior nota máxima foi a Sudeste ( 94,1 ), ao passo que a região que atingiu a menor nota máxima foi a Norte ( 87,3 ).
A mediana do Brasil como um todo foi 46,6, sendo a maior mediana obtida na região Sudeste ( 48,3 ), e a menor obtida na Centro-Oeste ( 43,6 ).
A nota mínima foi zero em todas as regiões, sem exceção.
Considerando -se as notas segundo Grande Região, observa -se que existe diferença estatisticamente significativa ao nível de 95 \% entre a maior média, obtida na região Sudeste ( 48,2 ), e as médias de todas as demais regiões.
O comportamento das notas dos estudantes de todo o Brasil pode ser observado no Gráfico 3.1 que apresenta um histograma com a distribuição das mesmas.
Essa é uma distribuição unimodal com a moda no intervalo ( 40 ; 50 ].
A Tabela 3.2 apresenta as estatísticas básicas da prova, desagregadas por Categoria Administrativa e por Organização Acadêmica.
Da população total de inscritos, 96.610 são de IES Privadas e 34.536 de IES Públicas.
Em relação à Organização Acadêmica, a maior participação foi obtida por estudantes de Universidades ( 79.420 ), seguido de Faculdades ( 33.151 ) e Centros Universitários ( 18.575 ).
A Categoria Administrativa de maior abstenção foi a Pública ( 15,4 \% ), e entre as Organizações Acadêmicas foi a dos Centros Universitários ( 17,7 \% ), ambos acima da média nacional de 14,7 \%.
A média das notas da prova como um todo foi 46,7.
Em relação à Categoria Administrativa, os estudantes das IES Públicas obtiveram média mais alta ( 49,1 ), e os das IES Privadas obtiveram média mais baixa ( 45,8 ), que a média nacional.
A diferença de 3,3 entre as médias das notas das IES Públicas e Privadas, é estatisticamente significativa.
Além disso, a diferença entre as médias das regiões Sudeste e Centro-Oeste ( 4,0 ), a maior e a menor média, é superior à diferença entre IES Públicas e Privadas ( 3,3 ), caracterizando uma maior diversidade regional do que administrativa.
No tocante à Organização Acadêmica, apenas as Faculdades ( 46,0 ) não obtiveram média mais alta que a nacional ( 46,7 ).
As demais médias foram : 47,0 nas Universidades e 46,9 nos Centros Universitários.
Os desvios padrões para as IES Públicas ( 15,4 ) e para as Universidades ( 15,3 ) foram superiores ao do Brasil como um todo ( 15,0 ), indicando uma dispersão um pouco maior das notas nesta Categoria Administrativa e nesta Organização Acadêmica.
Constata -se que existe diferença estatisticamente significativa ao nível de 95 \% nas médias das notas dos estudantes provenientes de Faculdades em relação aos de Centros Universitários e Universidades. 3.1.2 Estatísticas Básicas no Componente de Formação Geral A Tabela 3.3 apresenta as estatísticas básicas em relação ao componente da prova que avalia a Formação Geral dos estudantes concluintes.
Os estudantes de todo Brasil obtiveram desempenho médio de 49,2.
Quanto à variabilidade, o desvio padrão das notas dos estudantes do Brasil, como um todo, foi 17,2.
A maior média foi obtida na região Sudeste ( 50,9 ), e a menor, na região Norte ( 45,6 ).
As demais médias foram : 47,2 na região Nordeste ; 49,2 na região Sul ; e 47,9 na região Centro-Oeste.
Já o maior desvio padrão foi obtido na região Nordeste ( 17,5 ), e o menor, na região Sul ( 16,6 ).
Os demais desvios padrões foram : 17,3 na região Norte ; 17,1 na região Sudeste, e 17,0 na região Centro-Oeste.
A maior nota no Componente de Formação Geral da prova do ENADE/2014 foi 99,4, obtida por, pelo menos, um estudante na região Sudeste.
A menor nota máxima foi obtida na região Norte ( 96,8 ).
A mediana do Brasil, como um todo, foi 49,7, sendo a menor mediana encontrada na região Norte ( 45,7 ) ; e a maior encontrada na região Sudeste ( 51,5 ).
A nota mínima nesta parte foi zero em todas as regiões, sem exceção.
Considerando -se as notas segundo Grande Região, observa -se que existe diferença estatisticamente significativa entre a maior média das notas do Componente de Formação Geral, obtida na região Sudeste ( 50,9 ), e na região Norte ( 45,6 ), a menor média. 3.1.2 Estatísticas Básicas no Componente de Formação Geral A Tabela 3.3 apresenta as estatísticas básicas em relação ao componente da prova que avalia a Formação Geral dos estudantes concluintes.
Os estudantes de todo Brasil obtiveram desempenho médio de 49,2.
Quanto à variabilidade, o desvio padrão das notas dos estudantes do Brasil, como um todo, foi 17,2.
A maior média foi obtida na região Sudeste ( 50,9 ), e a menor, na região Norte ( 45,6 ).
As demais médias foram : 47,2 na região Nordeste ; 49,2 na região Sul ; e 47,9 na região Centro-Oeste.
Já o maior desvio padrão foi obtido na região Nordeste ( 17,5 ), e o menor, na região Sul ( 16,6 ).
Os demais desvios padrões foram : 17,3 na região Norte ; 17,1 na região Sudeste, e 17,0 na região Centro-Oeste.
A maior nota no Componente de Formação Geral da prova do ENADE/2014 foi 99,4, obtida por, pelo menos, um estudante na região Sudeste.
A menor nota máxima foi obtida na região Norte ( 96,8 ).
A mediana do Brasil, como um todo, foi 49,7, sendo a menor mediana encontrada na região Norte ( 45,7 ) ; e a maior encontrada na região Sudeste ( 51,5 ).
A nota mínima nesta parte foi zero em todas as regiões, sem exceção.
Considerando -se as notas segundo Grande Região, observa -se que existe diferença estatisticamente significativa entre a maior média das notas do Componente de Formação Geral, obtida na região Sudeste ( 50,9 ), e na região Norte ( 45,6 ), a menor média.
O Gráfico 3.2 propicia a avaliação do desempenho dos estudantes do Componente de Formação Geral, a partir do histograma da distribuição das notas correspondentes.
A distribuição é unimodal, com moda em ( 50 ; 60 ], enquanto na prova como um todo ( Gráfico 3.1 ) a moda foi alcançada dois intervalos abaixo ( 40 ; 50 ].
Destaca -se também o intervalo ( 40 ; 50 ] que possui valor muito próximo do modal.
Nota -se, ainda que, no Gráfico 3.2, as notas apresentam maior dispersão do que no Gráfico 3.1 ( distribuição das notas da prova ), confirmado pela comparação dos desvios padrões : 15,0 para a nota da prova como um todo e 17,2 para o Componente de Formação Geral.
Na Tabela 3.4 são apresentadas as informações referentes ao desempenho dos concluintes no Componente de Formação Geral, em diferentes agregações : Categoria Administrativa e Organização Acadêmica.
Observa -se que existe diferença estatisticamente significativa entre as médias dos tipos de Categoria Administrativa.
A maior média foi obtida por estudantes de IES Públicas ( 50,1 ), com uma diferença estatisticamente significativa da obtida por estudantes de IES Privadas ( 48,9 ).
Considerando -se o tipo de Organização Acadêmica, nota -se que não há diferença estatisticamente significativa entre as médias dos estudantes de Universidades ( 49,6 ) e de Centros Universitários ( 49,4 ) ; mas existe diferença entre a média desses e a dos de Faculdades ( 48,2 ). 3.1.3 Estatísticas Básicas do Componente de Conhecimento Específico A Tabela 3.5 apresenta as estatísticas básicas referentes ao Componente de Conhecimento Específico da área de Pedagogia ( Licenciatura ).
A média do desempenho dos estudantes do Brasil, como um todo, foi 45,8.
A maior média foi obtida na região Sudeste ( 47,3 ), e a menor, na região Centro-Oeste ( 42,9 ).
As demais médias foram : 43,9 na região Norte ; 45,0 na região Nordeste ; e 45,7 na região Sul.
Quanto à variabilidade das notas, o desvio padrão do Brasil, como um todo, foi 16,4, sendo o maior desvio padrão observado na região Sudeste ( 16,6 ), e o menor, nas regiões Norte e Centro-Oeste ( 16,0 ).
Os demais desvios foram : 16,3 na região Nordeste ; e 16,1 na região Sul.
A mediana das notas dos estudantes de todo o Brasil foi 45,8.
A maior mediana ocorreu na região Sudeste ( 47,3 ), e a menor, na região Centro-Oeste ( 42,4 ).
As demais medianas foram : 43,0 na região Norte ; 45,0 na região Nordeste ; e 45,7 na região Sul.
A nota máxima do Brasil, como um todo, foi 95,0, sendo obtida por, pelo menos, um estudante na região Sudeste.
As demais notas máximas foram : 87,8 na região Norte ; 92,8 na região Nordeste ; 91,3 na região Sul ; e 94,3 na região Centro-Oeste.
A nota mínima foi zero em todas as regiões, sem exceção.
Observa -se que existe diferença estatisticamente significativa da menor média das notas do Componente de Conhecimento Específico da região Centro-Oeste ( 42,9 ) em relação às demais regiões.
Assim como os Gráficos 3.1 e 3.2, o Gráfico 3.3, apresentado a seguir, permite uma avaliação do desempenho de concluintes de Pedagogia ( Licenciatura ), mas neste caso em relação ao Componente de Conhecimento Específico com um histograma da distribuição das notas correspondentes.
Dentre as três distribuições apresentadas, esta é a mais concentrada nas notas baixas.
Esta também é uma distribuição unimodal, e o grupo modal é o ( 40 ; 50 ], o mesmo grupo modal da prova como um todo e abaixo do grupo modal para a Formação Geral.
A Tabela 3.6 apresenta uma comparação dos resultados em relação à Categoria Administrativa e à Organização Acadêmica, agora levando em conta o desempenho dos estudantes no Componente de Conhecimento Específico da prova.
No que se refere à Organização Acadêmica, a maior média foi das Universidades ( 46,1 ), vindo a seguir a dos Centros Universitários ( 46,0 ) e, depois, a das Faculdades ( 45,2 ).
O maior desvio padrão, e acima do valor para o Brasil, como um todo, foi o das Universidades ( 16,8 ).
As Universidades também obtiveram a maior mediana ( 46,3 ) e a mesma nota máxima das Faculdades ( 94,3 ).
Por sua vez, os Centros Universitários obtiveram a maior nota máxima ( 95,0 ) e mediana igual a 46,0, enquanto as Faculdades obtiveram mediana igual a 45,0.
A nota mínima foi zero para todas as Organizações Acadêmicas, sem exceção.
Quanto à Categoria Administrativa, observa -se um comportamento semelhante àquele da parte de Formação Geral e à prova como um todo, ou seja, existe diferença estatisticamente significativa entre as médias das IES Públicas ( 48,7 ) e IES Privadas ( 44,8 ).
Como nos casos anteriores, a maior média foi obtida por estudantes de IES Públicas de ensino.
Observa -se que existem diferenças estatisticamente significativas ao nível de 95 \% no Componente de Conhecimento Específico entre as notas das Faculdades em relação aos Centros Universitários e Universidades. 3.2 ANÁLISE DAS QUESTÕES OBJETIVAS Esta seção apresenta estatísticas para o conjunto de questões objetivas e uma análise gráfica de questões selecionadas dos componentes de Formação Geral ( 3.2.1 ) e de Conhecimento Específico ( 3.2.2 ).
São também apresentadas e comparadas as médias das subpopulações caracterizadas por Grande Região. 3.2.1 Componente de Formação Geral A Tabela 3.7 apresenta as estatísticas básicas relativas às oito questões objetivas do componente da prova que abrange a Formação Geral dos estudantes.
A média do Brasil foi 48,3.
A menor média foi encontrada na região Norte ( 45,8 ), e a maior, na região Sudeste ( 49,6 ).
As demais médias foram : 47,2 na região Nordeste ; 48,0 na região Sul ; e 47,4 na região Centro-Oeste.
O desvio padrão do Brasil foi 20,2, sendo o maior desvio padrão encontrado na região Sudeste ( 20,4 ), e o menor, na região Sul ( 19,9 ).
Os demais desvios foram : 20,2 nas regiões Norte e Nordeste ; e 20,1 na região Centro-Oeste.
As medianas ( 50,0 ), as notas máximas ( 100,0 ) e as notas mínimas ( 0,0 ) foram iguais para todas as regiões, sem exceção.
A Tabela 3.8 apresenta o índice de facilidade e o índice de discriminação ( ponto bisserial ) para cada uma das questões objetivas do Componente de Formação Geral.
Quanto ao índice de facilidade, foram usadas as seguintes cores para diferenciar o nível de dificuldade da questão :  Azul para as questões classificadas com índice muito fácil ( \&gt ; =0,86 ), verde para as questões classificadas com índice fácil ( 0,61 a 0,85 ), amarelo para as questões classificadas com médio ( 0,41 a 0,60 ), vermelho para as questões classificadas com difícil ( 0,16 a 0,40 ) e roxo para as questões classificadas com muito difícil ( \&lt ; =0,15 ).
Já quanto ao índice de discriminação, foram usadas as seguintes cores para qualificar a questão :  As questões classificadas com índice fraco receberam a cor vermelho ( \&lt ; =0,19 ), as classificadas com médio receberam a cor amarelo ( 0,20 a 0,29 ), as classificadas com bom receberam a cor verde ( 0,30 a 0,39 ) e as classificadas com muito bom ( \&gt ; =0,40 ) receberam a cor azul.
As questões objetivas do Componente de Formação Geral, segundo o índice de facilidade, foram assim avaliadas : das oito questões, nenhuma teve o índice de facilidade classificado como muito fácil.
Duas questões foram tidas como fácil, por terem índice de acertos situado na faixa entre 0,61 e 0,85 ( de 61,0 \% a 85,0 \% de acertos ).
Três questões foram consideradas com índice de dificuldade médio, situando -se no intervalo entre 0,41 e 0,60 do índice de facilidade, ou seja, houve entre 41,0 \% e 60,0 \% de acertos, enquanto outras três questões foram classificadas na categoria difícil, situando -se no intervalo entre 0,16 e 0,40.
Por fim, nenhuma das questões apresentou menos de 15 \% de acertos, razão pela qual seria classificada como muito difícil.
Como já comentado, para análise das questões objetivas, relativas à Formação Geral, segundo o poder de discriminação, utilizou -se o índice de discriminação ( ponto bisserial ).
Nesta análise, para esse grupo de estudantes, as questões foram assim avaliadas : cinco das oito questões apresentaram índice maior ou igual a 0,40 e, assim, foram classificadas com índice muito bom ; três questões tiveram índice de discriminação bom, com valor entre 0,30 e 0,39 ; e nenhuma questão teve nível médio ou fraco de discriminação para esse grupo de estudantes.
O índice de facilidade variou de 0,25 a 0,80, e o de discriminação, de 0,34 a 0,49.
As cinco questões com índice de discriminação muito bom figuraram entre os diversos níveis de dificuldade desse conjunto : duas classificadas na categoria fácil ( questões 4 e 8 ) do índice de facilidade, duas, na categoria médio ( questões 3 e 6 ) e uma, na categoria difícil ( questão 7 ).
Em particular, a questão 7, junto da questão 6, foi a que apresentou o maior poder discriminatório, com índice 0,49, porém foi considerada difícil em termos de facilidade, com uma proporção de 0,39 acertos.
O máximo de acertos foi alcançado pela questão 8 com um índice de facilidade de 0,80.
A questão de número 1 apresentou um índice de facilidade de 0,48, ou seja, quase a metade dos estudantes conseguiu resolvê -la, dentro do universo de participantes.
Seu índice de discriminação foi bom ( 0,37 ).
Já a questão 5 obteve índice de discriminação bom, 0,34, porém seu índice de facilidade foi difícil ( 0,25 ).
O Gráfico 3.4, para exemplificar, analisa o comportamento da questão de número 7 de Formação Geral.
Trata -se de uma questão considerada difícil, com relação à facilidade e uma das duas que obtiveram o maior índice de discriminação dessa parte da prova.
Neste gráfico, cada uma das cinco curvas representa o percentual de respostas em determinada alternativa da questão, em função do número de acertos dos estudantes nessa parte da prova ( Formação Geral/Múltipla Escolha ), antes de possíveis eliminações pelo critério do ponto bisserial.
A curva em vermelho corresponde à alternativa E, a correta para esta questão.
Assim, observa -se que entre os estudantes com menor número de acertos, nessa parte do exame, a situação mais frequente foi a escolha de uma das alternativas incorretas : a alternativa C ( em preto ) ou D ( em roxo ).
Na medida em que o número de acertos aumenta, indicando desempenho melhor nessa parte da prova, aumenta concomitantemente a proporção de estudantes que selecionaram a alternativa correta E, atingindo 100 \% para os estudantes com 8 acertos.
Essa análise permite verificar como a questão discriminou os grupos de desempenho, justificando o alto índice obtido na questão.
Os gráficos relativos às demais questões de Formação Geral constam do Anexo I. 3.2.2 Componente de Conhecimento Específico A Tabela 3.9 apresenta as estatísticas básicas em relação às questões objetivas do Componente de Conhecimento Específico da prova de Pedagogia ( Licenciatura ) por Grande Região.
A média do Brasil deste componente foi de 49,6.
A menor média foi observada na região Centro-Oeste ( 46,5 ), e a maior, na região Sudeste ( 51,1 ).
O desvio padrão de todo o Brasil foi 18,3, sendo o menor desvio padrão encontrado na região Centro-Oeste ( 17,8 ), e o maior, na região Sudeste ( 18,5 ).
A mediana de todo o Brasil foi 50,0, a mesma encontrada nas regiões Nordeste, Sudeste e Sul.
Outras regiões apresentaram valores menores para a mediana : 45,5 nas regiões Norte e Centro-Oeste.
A nota máxima da prova foi 100,0, obtida nas questões objetivas do Componente de Conhecimento Específico, por, pelo menos, um estudante em todas as regiões, exceto na região Norte ( 95,5 ).
Em todas as regiões a nota mínima foi zero, sem exceção.
A Tabela 3.10 apresenta os índices de facilidade e discriminação ( ponto bisserial ) das questões objetivas do Componente de Conhecimento Específico para os estudantes de Pedagogia ( Licenciatura ).
Para facilitar a diferenciação das questões, foram usadas as mesmas cores da Tabela 3.8 para as diferentes classificações dos índices de facilidade e de discriminação.
Dentre as questões objetivas da parte da prova relativa ao Componente de Conhecimento Específico, nenhuma das questões foi anulada pela Comissão.
Desse modo, a classificação, quanto ao índice de facilidade, foi estabelecida com base em todas as 27 questões.
A partir dos índices obtidos, pode -se concluir que um pouco menos da metade das questões objetivas da prova foi considerado, pelo menos, difícil : das 27 questões, 11 classificadas como difícil ( a classificação modal ) e duas como muito difícil.
Sete questões foram classificadas como fácil, e outras sete consideradas como médio.
Não houve questão classificada como muito fácil.
Já quanto aos índices de discriminação das questões objetivas do Componente de Conhecimento Específico da prova, tem -se como resultado a seguinte classificação : oito das 27 questões foram consideradas como boas, enquanto 11 delas tiveram índice de discriminação muito bom.
Assim, para a maioria das questões – 19 em 27 – os índices de discriminação foram bom ou muito bom.
Dentre as demais, três delas foram classificadas como médio e outras cinco como fraco, sendo oito, por conseguinte, a quantidade de questões nos dois patamares mais baixos de discriminação.
Constata -se, assim, que a prova – no que se refere ao Componente de Conhecimento Específico – possuía razoável capacidade de discriminar entre aqueles que dominam ou não o conteúdo.
Dentre as questões que alcançaram os maiores índices de discriminação, 11 delas, as de números 11, 12, 13, 16, 18, 20, 22, 26, 31, 34 e 35, foram classificadas com índice de discriminação muito bom, situando -se no intervalo de 0,40 a 0,49 do índice, quatro delas ( questões 12, 16, 22 e 26 ) foram classificadas na categoria fácil, quanto ao índice de facilidade, cinco delas ( questões 11, 13, 18, 20 e 31 ), na categoria médio, e as outras duas ( questões 34 e 35 ), na categoria difícil.
A questão de número 17 foi a mais difícil dentre as 27 questões específicas, com baixo índice de facilidade, apenas 12,0 \% de acertos.
Essa questão apresentou poder discriminatório igualmente baixo, 0,04, o que comprova ter sido esta questão a mais difícil para os estudantes.
Destaca -se, também, a questão 33, com índice de facilidade 0,14, o que, em termos percentuais, corresponde a 14,0 \% de estudantes que responderam acertadamente.
Já, 0,05 foi o seu índice de discriminação.
Tais questões foram, portanto, pelo critério ponto bisserial, consideradas inadequadas.
Por isso, as questões 17 e 33 foram eliminadas do cômputo da nota final.
Além destas duas, as demais questões com índice de discriminação fraco, questões 14, 25 e 30 tampouco foram utilizadas no cômputo final das notas, num total de cinco questões eliminadas.
A título de exemplo das análises do comportamento das questões objetivas, o Gráfico 3.5 analisa a questão 26 do Componente de Conhecimento Específico.
Esta questão foi classificada como fácil.
Apresentou índice de facilidade 0,67, ou seja, 67,0 \% dos estudantes assinalaram acertadamente a opção E, correspondente ao gabarito.
Seu índice de discriminação foi igual a 0,49, classificado como muito bom, também sendo esta questão a que apresentou o maior índice discriminatório.
Neste gráfico, cada uma das cinco curvas representa o percentual de respostas em determinada alternativa da questão 26, em função do número de acertos dos estudantes nessa parte da prova, antes de possíveis eliminações de questões pelo critério do ponto bisserial.
A alternativa correta E, representada no gráfico pela curva em vermelho, foi escolhida em maiores proporções pelos estudantes com desempenho melhor nessa parte da prova.
Já as alternativas incorretas, também denominadas distratores, foram selecionadas, principalmente, por aqueles com notas mais baixas.
Neste caso a soma não é sempre 100 \% por causa das questões não respondidas ou com mais de uma opção marcada.
Aqueles com nota zero, na sua quase totalidade deixaram esta questão em branco ou marcaram mais de uma alternativa, comportamento considerado inválido.
A proporção de estudantes que selecionou a resposta correta E aumenta gradativamente, chegando a atingir 100 \% para 23 acertos ou mais, enquanto a proporção dos que escolheram alternativas incorretas decai, a partir de dois acertos, como função do número de acertos nessa parte da prova.
Os gráficos relativos às demais questões do Conhecimento Específico constam do Anexo I. 3.3 ANÁLISE DAS QUESTÕES DISCURSIVAS Esta seção apresenta estatísticas selecionadas e histogramas das Questões Discursivas de Formação Geral ( 3.3.1 ) e Conhecimento Específico ( 3.3.2 ).
São também apresentadas e comparadas as médias de subpopulações caracterizadas por Grande Região. 3.3.1 Componente de Formação Geral As análises dos resultados de desempenho dos estudantes de Pedagogia ( Licenciatura ) nas duas questões discursivas relativas à Formação Geral, considerando -se as notas de conteúdo e de desempenho linguístico, encontram -se na Tabela 3.11 e no Gráfico 3.6.
Na Tabela 3.11, observa -se que a nota média nesse conjunto de questões foi próxima da obtida nas objetivas.
Os estudantes, de todo o Brasil, obtiveram, em Formação Geral, média 48,3 nas questões objetivas e 50,6 nas questões discursivas.
No entanto, pode -se notar um aumento do desvio padrão de 20,2 nas questões objetivas do Componente de Formação Geral dos estudantes de todo o Brasil, para A mediana de todo o Brasil, neste componente, foi 54,0, na região Sudeste a mediana foi maior, 56,5.
Nas demais regiões foi menor : Norte ( 48,5 ), Nordeste ( 51,0 ), Sul ( 53,5 ) e Centro-Oeste ( 52,0 ).
A nota máxima ( 98,5 ) foi obtida por pelo menos um estudante na região Sudeste, enquanto nas demais regiões a nota máxima foi : 97,5 nas regiões Norte e Centro-Oeste ; e 98,0 nas regiões Nordeste e Sul.
A nota mínima ( 0,0 ) foi a mesma em todas as regiões do Brasil, sem exceção.
O Gráfico 3.6 representa a distribuição das notas nas questões discursivas do Componente de Formação Geral.
A moda desta distribuição ocorre no intervalo ( 50 ; 60 ].
Destacam -se, também, o intervalo ( 60 ; 70 ], com frequência bem próxima a do valor modal, e o intervalo [ 0 ; 10 ] com distribuição próxima a 10 \% do total de notas, sendo que no intervalo [ 0 ; 10 ] incluem -se além da nota zero, a frequência de estudantes que deixaram este tipo de questão em branco. 23,3 nas questões discursivas do mesmo componente.
A maior média foi obtida na região Sudeste ( 53,0 ), e a menor, na região Norte ( 45,4 ).
A mediana de todo o Brasil, neste componente, foi 54,0, na região Sudeste a mediana foi maior, 56,5.
Nas demais regiões foi menor : Norte ( 48,5 ), Nordeste ( 51,0 ), Sul ( 53,5 ) e Centro-Oeste ( 52,0 ).
A nota máxima ( 98,5 ) foi obtida por pelo menos um estudante na região Sudeste, enquanto nas demais regiões a nota máxima foi : 97,5 nas regiões Norte e Centro-Oeste ; e 98,0 nas regiões Nordeste e Sul.
A nota mínima ( 0,0 ) foi a mesma em todas as regiões do Brasil, sem exceção.
O Gráfico 3.6 representa a distribuição das notas nas questões discursivas do Componente de Formação Geral.
A moda desta distribuição ocorre no intervalo ( 50 ; 60 ].
Destacam -se, também, o intervalo ( 60 ; 70 ], com frequência bem próxima a do valor modal, e o intervalo [ 0 ; 10 ] com distribuição próxima a 10 \% do total de notas, sendo que no intervalo [ 0 ; 10 ] incluem -se além da nota zero, a frequência de estudantes que deixaram este tipo de questão em branco.
Na sequência, os resultados verificados para cada uma das questões discursivas de Formação Geral serão apresentados, estabelecendo -se relações com os conteúdos abordados em cada uma delas.
Os comentários da Banca de docentes corretores, a respeito do observado na correção das respostas dos estudantes, suas impressões e conclusões serão apresentados junto à análise de cada questão.
Cumpre esclarecer que, tendo em vista que as questões discursivas de Formação Geral são padronizadas, ou seja, constam de todas as provas, os comentários da Banca são os mesmos para todas as carreiras acadêmicas, sendo direcionados a todos os estudantes que participaram do ENADE/2014.
A seguir, serão analisados os desempenhos da Área de Pedagogia ( Licenciatura ) nas duas questões discursivas de Formação Geral do ENADE/2014, comparando -se os resultados obtidos com comentários para cada questão. 3.3.1.1 Análise de Conteúdo da Questão Discursiva 1 do Componente de Formação Geral Os dados de Pedagogia ( Licenciatura ), obtidos a partir das respostas à questão 1, encontram -se na Tabela 3.12 e no Gráfico 3.7.
Nessa questão – de pior desempenho dentre as duas de Formação Geral – os estudantes, de todo o Brasil, tiveram média 47,5.
A maior média para a questão 1 foi obtida na região Sudeste ( 50,7 ), e a menor, na região Norte ( 39,9 ).
Quanto à variabilidade das notas, o desvio padrão de todo o Brasil foi 31,0.
O menor desvio padrão foi obtido na região Sul ( 30,2 ), e o maior desvio padrão foi obtido na região Nordeste ( 31,8 ).
As medianas das regiões Sudeste, Sul, Centro-Oeste e do Brasil, como um todo, foram iguais ( 50,0 ).
Nas regiões Norte e Nordeste, as medianas foram 40,0 e 45,0 ; respectivamente.
As notas máximas e mínimas da questão discursiva 1 foram as mesmas para todas as regiões do Brasil, respectivamente, 100,0 e 0,0, sem exceção.
O Gráfico 3.7 mostra a distribuição das notas na questão discursiva 1 do Componente de Formação Geral.
Observa -se que a maior frequência corresponde aos estudantes que obtiveram nota dentro do intervalo ( 70 ; 80 ], correspondendo à moda da distribuição.
Destacam -se, também, os estudantes que deixaram a questão em branco, cerca de 10 \% do total, e aqueles cuja nota ficou situada no intervalo ( 40 ; 50 ], quase 15 \% do total. 3.3.1.2 Comentários sobre a correção de Conteúdo das respostas à Questão Discursiva 1 O enunciado apresentado da Questão 1 era claro e foi redigido em linguagem acessível.
A questão abordava o tema da mobilidade urbana, relacionando -a aos ideais de desenvolvimento urbano sustentável estabelecidos pela ONU, do ponto de vista ecológico, cultural, político, institucional, social e econômico.
O enunciado apresentava uma tabela, sem referência à fonte de publicação, que detalhava percentualmente as modalidades de deslocamento da população urbana brasileira, com base em duas distinções : motorizado versus não motorizado ; coletivo versus individual.
O comando da questão solicitava, então, que o estudante desenvolvesse um texto dissertativo composto de duas partes : análise das consequências do transporte motorizado para o desenvolvimento sustentável e elaboração de duas propostas de intervenção pautadas pelo incentivo ao uso da bicicleta, meio de transporte não motorizado e, portanto, não poluente.
A temática dos efeitos da industrialização, com foco na produção de automóveis, e consequências para o meio ambiente é atual.
Além disso, a questão provoca a reflexão sobre a necessidade de adequação das cidades ao número de habitantes, no contexto da utilização dos meios de transportes de massa : metroviário, ferroviário, por ônibus.
Esperava -se, assim, que o concluinte recorresse a diferentes argumentos para analisar as consequências do transporte motorizado, explicitados pelo texto motivador, e propusesse duas ações de intervenção, as quais considerasse relevantes para o incremento ao uso de bicicleta, que, preferencialmente, deveriam manter coerência com a argumentação desenvolvida na primeira parte da resposta.
O conflito entre a busca pelo transporte individual e o coletivo é inerente a todas as grandes cidades do mundo.
No Brasil, esse conflito emerge com maior gravidade, como efeito do aumento da quantidade de veículos automotores ( automóveis e motocicletas ) e da baixa qualidade dos transportes coletivos.
Tal associação incentiva o uso de meios de transporte individual em detrimento da utilização dos coletivos.
O padrão de respostas era tecnicamente adequado ao enunciado e amplo o suficiente para abranger as variações de respostas apresentadas.
Deve ser ressaltado que a proteção ao ambiente ecologicamente equilibrado é tema conhecido dos concluintes dos cursos universitários, que deveriam possuir informações suficientes para desenvolver uma resposta abordando pelo menos dois dos aspectos listados no padrão de respostas como possibilidades de consequências ( item ‘ a ’ ) e duas propostas de ação de intervenção aceitáveis ( item ‘ b ’ ).
Segundo o padrão, concluintes poderiam desenvolver, quanto ao item ‘ a ’, os seguintes pontos : aumento da emissão de poluentes atmosféricos ; aumento da emissão de gases de efeito estufa ( CO2- \textbf{dióxido} de \textbf{carbono}, CO - monóxido de \textbf{carbono}, O3 - Ozônio ) ; aumento da poluição visual e sonora ; aumento da temperatura local e global ; aumento do consumo de combustíveis ; aumento de problemas de saúde ( cardíacos, respiratórios, dermatológicos ) ; aumento da frota de veículos, promovendo congestionamentos urbanos ; diminuição de áreas verdes ; desmatamento ; aumento das áreas de impermeabilização, resultando em enchentes, diminuição da infiltração da água e recarga de lençóis freáticos ; elevação dos custos de manutenção das cidades ( metroferrovias, rodovias, tratamento de água, limpeza da cidade, etc. ) ; necessidade de ampliação de vias trafegáveis ; necessidade de ampliação de áreas de estacionamento.
Em relação ao item ‘ b ’, os estudantes poderiam considerar os seguintes aspectos : construção de vias exclusivas para bicicletas ( ciclovias e ciclofaixas ) ; proposição de formas de integração entre o transporte por bicicletas, o metroviário e os ônibus coletivos, a fim de garantir segurança e conforto em momentos de adversidades climáticas e relevo acidentado ; pontos de aluguel e/ou empréstimos de bicicletas ; construção de bicicletários ; investimentos na segurança pública ; políticas de incentivo ao uso de bicicletas ( educação ambiental, qualidade de vida, saúde, propaganda ) ; implementação de políticas de crédito e de redução de custos das bicicletas.
Para o item ‘ a ’, foi encontrado um número expressivo de respostas adequadas.
As consequências mais citadas estavam relacionadas com a emissão de poluentes e o aumento dos engarrafamentos.
Houve críticas frequentes à qualidade do transporte público e à necessidade de locomoção rápida diante da distância das residências aos locais de trabalho.
Em tom de reivindicação, muitos estudantes defenderam solução para os problemas do transporte coletivo, quer metroviário, ferroviário ou por ônibus.
O uso da bicicleta também foi considerado por muitos como um meio de transporte benéfico à saúde, mas alguns apontaram a impossibilidade de sua utilização pelos idosos e deficientes.
No concernente ao item ‘ b ’, de modo quase uniforme citou -se a necessidade de ciclovias e ciclofaixas, bem como de se garantir segurança para os ciclistas.
A necessidade de manutenção das ciclovias e ciclofaixas existentes em algumas cidades foram também lembradas.
Com menos frequência também houve a indicação de maior conscientização da população, inclusive pela mídia, e pelo incentivo à aquisição desses veículos não motorizados, por meio da diminuição de impostos para redução do custo de aquisição.
Houve quem sugerisse, inclusive, uma política de doação de bicicletas.
Alguns mencionaram a necessidade de se assegurar integração das ciclovias com os demais modais, com a construção de bicicletários seguros diante das grandes distâncias a percorrer, notadamente nas metrópoles.
Também com base na dificuldade de o veículo não motorizado vencer as longas distâncias que o trabalhador deve percorrer nos centros urbanos para chegar ao local onde exerce a sua atividade produtiva, alguns estudantes argumentaram contra a utilização da bicicleta como forma de resolução da poluição ambiental.
O principal equívoco cometido na sugestão de ações ( item ‘ b ’ ) foi o de mencionar a necessidade de melhoria do meio ambiente dissociado do uso de bicicletas, como se requeria no comando da questão.
Nesse sentido, foram encontradas respostas que indicavam a \textbf{caminhada} como opção.
Foram consideradas respostas fracas ( notas de zero a 30 ) aquelas meramente opinativas, ou que apenas transcreveram elementos do enunciado, sem qualquer acréscimo ou apreciação crítica.
As respostas medianas ( notas entre 35 e 70 ) foram as que não apresentavam algum desenvolvimento para os tópicos citados ( duas consequências e duas ações de intervenção ) ou por falta de alguns deles.
Já as respostas boas ( notas de 75 a 100 ) continham os tópicos considerados pelo padrão de respostas, com algum desenvolvimento, pelo menos, regular de argumentação.
As respostas dos estudantes revelaram um adequado conhecimento quanto aos efeitos do transporte motorizado para o ambiente e que a temática da ecologia parece ser tema bem compreendido, aparecendo nas respostas como uma preocupação para as futuras gerações.
Nesse sentido, algumas políticas públicas foram mencionadas apropriadamente para assegurar o desenvolvimento sustentável pelos respondentes, tais como : educação, habitação, saneamento, saúde, transporte e segurança.
No entanto, de modo geral, as respostas demostraram que os estudantes têm muita dificuldade na expressão escrita do pensamento, como se pode constatar pela avaliação do desempenho linguístico que ficou a cargo de uma banca específica, formada por profissionais da área de Língua Portuguesa. 3.3.1.3 Análise de Conteúdo da Questão Discursiva 2 do Componente de Formação Geral A Tabela 3.13 mostra que o desempenho médio dos estudantes na questão discursiva 2 ( média 51,1 ) foi superior ao obtido na questão discursiva 1 ( média 47,5 ).
A região Sudeste foi aquela cuja média, nessa questão, foi maior ( 53,0 ), e a de menor média foi a região Norte ( 48,2 ).
Quanto à variabilidade das notas, o desvio padrão de todo o Brasil foi 30,1, inferior ao obtido na questão discursiva 1 ( 31,0 ).
O maior desvio nessa questão foi obtido na região Norte ( 32,0 ), enquanto o menor foi obtido na região Sul ( 29,0 ).
A mediana de todo o Brasil foi 60,0, a mesma da região Sudeste, enquanto nas demais regiões a mediana foi 50,0.
As notas máximas ( 100,0 ) foram as mesmas em todas as regiões do Brasil, assim como as notas mínimas ( 0,0 ), sem exceção.
O Gráfico 3.8 mostra a distribuição das notas na questão discursiva 2 do Componente de Formação Geral.
Observa -se que a maior frequência corresponde aos estudantes que obtiveram nota situada no intervalo ( 40 ; 50 ]. 3.3.1.4 Comentários sobre a correção de Conteúdo das respostas à Questão Discursiva 2 A questão 2 abordava o tema da violência urbana a partir de um texto motivador recolhido em um site de um jornal brasileiro.
O fragmento relatava a prisão de três jovens por agressão a um outro jovem, menor de idade, guardador de carros.
A notícia trazia informações sobre o estado do jovem agredido e sobre as possíveis penas a serem aplicadas aos agressores.
O enunciado solicitava, então, que o estudante desenvolvesse um texto dissertativo composto de duas partes : análise de duas causas do tipo de violência descrito no texto e apresentação de dois fatores que contribuiriam para evitar a agressão relatada.
Como o texto motivador reduzia -se ao relato de um episódio de violência urbana, o enunciado não oferecia argumentos que \textbf{ajudassem} o estudante a elaborar uma análise das possíveis causas do fato.
As propostas a serem apresentadas, no sentido de se evitar em situações como a relatada, deveriam manter coerência com a análise das causas, desenvolvida na primeira parte da resposta.
O enunciado era claro e a matéria permitia identificar tratar -se de um conflito tipicamente urbano no qual se envolveram jovens de baixa renda em disputa por um território onde exerciam trabalho informal.
Sendo essa uma das hipóteses albergada sobre o amplo tema da violência urbana no qual se desenvolve o relato.
A temática da questão é bastante conhecida, não somente por debates no âmbito universitário, especialmente na área de Ciências Sociais, bem como pela ênfase que é dada pela mídia onde assume contornos, muitas vezes, maiores do que o problema em si.
Assim, por tratar -se de tema bastante presente e debatido no cotidiano das pessoas, a questão pode ser considerada de baixa complexidade.
O padrão de resposta apresentava uma gama de possibilidades de causas da violência aceitas como corretas para o item ‘ a ’, tais como : problemas relacionados à educação ( baixa escolaridade, evasão escolar, qualidade da educação, distanciamento entre escola e a realidade social, tempo de permanência na escola ) ; desigualdades socioculturais ( gênero, etnia, economia, etc. ) ; desemprego e falta de qualificação profissional ; precariedade da segurança pública ; uso de drogas ; desvalorização da vida humana ; banalização da violência ; sensação de impunidade ; ausência de políticas sociais ; degradação da vida urbana ; desconhecimento ou desrespeito aos direitos humanos e constitucionais ; desestruturação familiar ; desvalorização de princípios éticos e morais.
No concernente ao item ‘ b ’, o padrão também previa a possibilidade de apresentação de diversos fatores que podem contribuir para evitar a violência : políticas de segurança mais efetivas ; políticas públicas de melhoria das condições socioeconômicas ; maior consciência cidadã e respeito à vida ; melhor distribuição de renda ; melhoria da educação ( aumento da escolaridade, redução da evasão escolar, qualidade da educação, aproximação entre a escola e a realidade social, aumento do tempo de permanência na escola ) ; aumento da oferta de emprego e melhoria da qualificação profissional ; medidas preventivas ao uso de drogas ; maior eficácia do sistema judiciário ; revisão da legislação penal ; valorização de princípios éticos, morais e familiares.
O enunciado da questão permitia que os estudantes, em suas respostas, pautassem o problema por diversos ângulos.
Dessa forma, o padrão buscou indicar diversas possibilidades a serem identificadas pelo corretor na leitura da resposta do estudante.
As notas foram atribuídas considerando o desenvolvimento da resposta e a indicação de ao menos dois dos itens exigidos pelo comando da questão nos itens ‘ a ’ ( duas causas ) e ‘ b ’ ( dois fatores ), dentre os considerados pelo padrão de respostas.
No entanto, o padrão exigia que as respostas estivessem afinadas com o perfil profissional previsto no Art. 3º da Portaria nº 255, de 02 de junho de 2014, onde se lê : Art. 3º No componente de Formação Geral serão considerados os seguintes elementos integrantes do perfil profissional : atitude ética ; comprometimento social ; compreensão de temas que transcendam ao ambiente próprio de sua formação, relevantes para a realidade social ; espírito científico, humanístico e reflexivo ; capacidade de análise crítica e integradora da realidade ; e aptidão para socializar conhecimentos em vários contextos e públicos diferenciados.
Nesse sentido, teses repressivas não conformadas com o texto constitucional e amparados pelos direitos humanos, não foram consideradas.
Questões como a majoração das penas, instituição de penas mais duras e criminalização da conduta dos menores desviados, só foram consideradas no caso de serem tratadas como menção a um debate que está presente, nunca como medida a ser aplicada de forma arbitrária e inconstitucional.
Também não foram previstas visões mais autoritárias e repressoras como a defesa de proibição do trabalho informal, como é o caso dos “ flanelinhas ”, e a defesa de não fornecer dinheiro aos “ guardadores ” de carro, identificado como esmolas por alguns.
Não foram consideradas respostas que meramente reproduziam o enunciado sem qualquer reflexão ou análise.
Aquelas que utilizaram os elementos do enunciado, mesmo que apenas transcritos, mas desenvolveram algum raciocínio pessoal ou indicaram causas e/ou fatores foram valoradas.
As causas apontadas com maior frequência para a violência foram : a ausência de educação e o desemprego, com soluções correlatas de educação profissional e de qualidade e abertura de postos de trabalho.
A questão das drogas também foi bastante mencionada, e a solução correlata mais indicada foi o acolhimento de usuários de drogas.
Outro tema bastante citado foi a dissolução das famílias e dos valores familiares.
Tendo em vista que o texto falava de moradores de rua, a ausência de moradia apareceu em muitas respostas, bem como a disputa territorial e a necessidade de sobrevivência.
Nesse sentido, as propostas apontavam para a necessidade de ajuda da área de Serviço Social das prefeituras e atendimento psicológico gratuito.
Numa ótica mais repressora, surgiram críticas quanto à leniência das leis, acompanhadas da indicação da necessidade de atuação mais rigorosa do Poder Judiciário, como também da aplicação de penas maiores, negativa de esmolas, proibição do trabalho informal ( guardador de carros ), dentre outras.
Algumas respostas indicaram como causa o modelo capitalista, o consumismo e a ganância.
Outras indicaram o trabalho infantil e a correlata necessidade de aplicação do ECA ( Estatuto da Criança e do Adolescente ).
Em relação ao tema violência, inúmeras respostas apresentaram visões de intolerância quanto aos menores infratores, postulando a diminuição da idade para efeitos de criminalização, bem como defendendo a majoração de penas.
Essa perspectiva revela uma cisão social, de certa forma estimulada pelo tipo de divulgação da mídia quanto aos temas em foco, o que, sem dúvida, ficou evidenciado em algumas respostas.
A existência de inúmeras respostas pugnando pela maior repressão estatal e pelo aumento das penas, acrescido da busca pela diminuição da idade penal, revela que muitos estudantes não foram atingidos por uma formação superior que se exige humanista, ética e comprometida socialmente, que busque os fundamentos para resolver o problema da violência por meio de práticas democráticas e de inclusão social.
As respostas desse tipo também revelaram um alto grau de intolerância para com os moradores de rua, associando -os muitas vezes ao tráfico de drogas, à exploração indevida do espaço público, bem como a crimes como extorsão dos motoristas quando pedem dinheiro para guardar os veículos.
Nessa trilha, verifica -se lacuna na discussão de temas importantes para a formação de profissionais de nível superior tais como : sociodiversidade, multiculturalismo e violência ; tolerância/intolerância ; inclusão/exclusão.
Para aqueles que responderam em conformidade com o padrão de resposta verificou -se uma plena formação social, conclamando pelo auxílio não só do Estado, mas também da Sociedade no amparo daqueles que não têm as mínimas condições de sobrevivência e que estão nas ruas. 3.3.1.5 Análise de Língua Portuguesa das Questões Discursivas do Componente de Formação Geral Os dados de Pedagogia ( Licenciatura ), obtidos a partir das respostas às questões discursivas do Componente de Formação Geral, no que tange à Língua Portuguesa, encontram -se na Tabela 3.14 e no Gráfico 3.9.
Nesse aspecto, os estudantes, de todo o Brasil, tiveram média 55,7.
A maior média com respeito à Língua Portuguesa foi obtida na região Sudeste ( 57,5 ), e a menor, na região Norte ( 50,9 ).
Quanto à variabilidade das notas, o desvio padrão de todo o Brasil foi 20,7.
O menor desvio padrão foi obtido na região Sul ( 18,7 ) e o maior desvio padrão foi obtido na região Nordeste ( 22,8 ).
A mediana das notas de Língua Portuguesa foi 60,0 para quase todas as regiões do Brasil, menos para as regiões Norte ( 57,5 ) e Sudeste ( 62,5 ).
A nota máxima para todo o Brasil foi de 95,0, com, pelo menos, um estudante tirando essa nota nas regiões Nordeste e Sudeste, enquanto nas demais regiões a nota máxima foi : 90,0 na região Norte ; e 92,5 nas regiões Sul e Centro-Oeste.
Além disso, a nota mínima foi zero em todas as regiões do país, sem exceção.
O Gráfico 3.9 mostra a distribuição das notas de Língua Portuguesa do Componente de Formação Geral.
Observa -se que a maior frequência ( cerca de 30 \% ) corresponde aos estudantes que obtiveram nota no intervalo ( 60 ; 70 ]. 3.3.1.6 Comentários sobre a correção das respostas de Formação Geral com respeito à Língua Portuguesa Ao encaminhar as questões 1 e 2 na direção da produção de um texto dissertativo, esperava -se que o estudante utilizasse seus conhecimentos sobre o assunto e estruturasse seus textos de acordo com as exigências do registro formal próprio dessa situação comunicativa.
Essa configuração determina exigências quanto : à adequação da seleção vocabular, ao desenvolvimento do conteúdo, à estruturação sintática dos períodos, à organização lógica das ideias, à utilização de procedimentos de encadeamento textual e referenciação, à obediência às exigências morfossintáticas próprias da modalidade escrita da norma-padrão, ao respeito às regras ortográficas e às regras de acentuação gráfica.
O padrão de respostas utilizado na avaliação considerou os aspectos relevantes ao bom desempenho linguístico como competências distintas, de modo a permitir um mapeamento detalhado do domínio dos recursos disponíveis na Língua Portuguesa para a comunicação escrita formal.
Com base nesse objetivo, foram avaliados os seguintes aspectos : a ) Estruturação textual condizente com o gênero solicitado e o modo de organização textual expositivo adequado ao gênero – essa competência envolve : a estruturação sintática condizente com o padrão da modalidade escrita formal da língua portuguesa, de modo a garantir a clareza necessária ; a distribuição do conteúdo do texto em parágrafos, de modo a garantir a sua organização temática ; a utilização de operadores discursivos que contribuam para a progressão temática do texto, estabelecendo relações lógicas entre as ideias apresentadas, tanto do ponto de vista intrafrasal, como do interfrasal ; a utilização de procedimentos de referenciação lexical e pronominal que permitam a retomada de referentes textuais ; o respeito às regras de pontuação como fator de estruturação do período.
Espera -se, portanto, que o estudante recorra a procedimentos linguístico-discursivos para organizar seu texto, permitindo o encadeamento lógico entre suas partes de forma a garantir a progressão e a coerência textuais.
Isso significa que os seguintes procedimentos foram penalizados, de acordo com o padrão de respostas proposto :  a estrutura lógico-gramatical do texto fica comprometida pela elaboração de frases fragmentadas ;  sequência justaposta de ideias sem encaixamentos sintáticos, reproduzindo hábitos da oralidade ;  elaboração de frase com apenas oração subordinada, sem oração principal ;  emprego equivocado do conector ( preposição, conjunção, pronome relativo, alguns advérbios e locuções adverbiais ) que não estabeleça relação lógica entre dois trechos do texto e prejudique a compreensão da mensagem ;  emprego do pronome relativo sem a preposição, quando obrigatória ;  repetição ou substituição inadequada de palavras sem utilização dos recursos oferecidos pela língua ( pronome, advérbio, artigo, sinônimo ) ;  emprego inadequado dos pronomes relativos “ cujo(a) ” e “ onde ” ;  utilização inadequada dos sinais de pontuação que comprometa a clareza textual. b ) Respeito às convenções ortográficas da norma-padrão da Língua Portuguesa – essa competência envolve o domínio das regras de acentuação gráfica e da grafia padrão das palavras ( com ausência de abreviaturas próprias da linguagem da internet ), de acordo com as convenções estabelecidas pela legislação em vigor e consubstanciadas no Vocabulário Ortográfico da Língua Portuguesa, editado pela Academia Brasileira de Letras ( com aceitação da legislação anterior, no caso das regras relativas ao uso do hífen e da acentuação gráfica ).
Espera -se que o participante :  grafe corretamente as palavras ;  respeite as regras de acentuação gráfica ;  empregue maiúsculas em início de frase, em nomes próprios de pessoas, lugares ou instituições ;  evite abreviações como p/, vc, tb, pq, tá, né, usadas muitas vezes em escrita informal e na internet ;  obedeça às regras de separação de sílabas no final da linha. c ) Domínio dos diferentes aspectos morfossintáticos próprios da modalidade escrita formal da norma-padrão da Língua Portuguesa – essa competência envolve : a concordância nominal, a concordância verbal, a regência nominal, a regência verbal, a flexão nominal, a flexão verbal, a correlação entre os tempos verbais, a colocação pronominal e a utilização de sinais de pontuação que contribuam para a organização lógica da frase e do texto.
Espera -se que o participante :  flexione o verbo para estabelecer concordância de número com o sujeito da frase ;  flexione o artigo, o adjetivo e o pronome para concordar em número e em gênero com o substantivo a que se referem ;  observe a regência nominal e a verbal, utilizando a preposição adequada depois de um substantivo, um verbo ou um adjetivo ;  empregue adequadamente o acento grave indicador de crase nos casos em que se fizer necessário ;  obedeça às regras de colocação pronominal ( próclise e ênclise ), distintas dos hábitos da oralidade ou da escrita informal ; d ) Seleção vocabular adequada à modalidade escrita formal da Língua Portuguesa, exigida pela situação comunicativa – essa competência envolve a precisão na utilização do vocabulário relacionado à temática solicitada pela questão ; a ausência de marcas da oralidade, como termos de sentido de muito genérico ( “ coisa ”, “ negócio ”, “ você ” ) e termos de registros mais informais ( como gírias, jargões, frases feitas, ditados populares, termos regionais ).
Assim, espera -se que o participante respeite a adequação vocabular não utilizando gírias ou expressões coloquiais, evite repetição desnecessária de palavras e utilize um vocabulário mais formal, adequado ao texto de caráter dissertativo.
A escolha dessas competências para subsidiar o processo de avaliação apoia -se na concepção de que, no desempenho dos graduandos, a modalidade escrita tem apresentado uma intensa simplificação, originada no padrão da modalidade oral da Língua Portuguesa.
No caso do texto de base dissertativa, inscrito em um registro formal, a distância entre as duas modalidades é ainda maior, o que provoca situações de hipercorreção ( desvios provocados pela incorporação indevida de uma regra da norma-padrão ) e de truncamentos sintáticos ( estruturas frasais incompreensíveis devido à complexidade sintática própria da modalidade escrita ).
Observam -se, então, os seguintes aspectos que marcam essa distinção entre as duas modalidades, devido à excessiva simplificação da modalidade falada : a ) redução drástica de estruturas subordinadas, compensada pelo aumento na frequência de estruturas coordenadas e absolutas, por um lado, ou pela elaboração de estruturas truncadas pelo excesso de ideias sem a devida conexão subordinativa ; b ) redução no uso de conectores para expressar relações lógicas essenciais à construção do texto, substituídas pela exigência de inferência por parte do interlocutor para suprir a sua ausência ; c ) redução cada vez maior do uso do subjuntivo, ao lado da ampliação do uso do indicativo combinado a estruturas frasais coordenadas ou absolutas ; d ) empobrecimento do processo de referenciação, com a repetição exaustiva de pronomes ou nomes ; e ) simplificação extrema da marcação da categoria tempo na morfologia verbal ; f ) falta de domínio de vocabulário mais abstrato e de maior complexidade, essencial ao desenvolvimento do processo dissertativo ; g ) redução drástica no emprego da acentuação gráfica, processo intensificado pela divulgação imprecisa das mudanças promovidas pelo último acordo ortográfico.
Os aspectos macroestruturais da elaboração do texto não foram avaliados neste processo, para não penalizar duplamente os estudantes, já que a banca de formação geral, composta por profissionais de diferentes áreas do conhecimento, encarregou -se da avaliação do conteúdo desenvolvido nas questões.
São eles : progressão temática, coerência na relação com os conhecimentos de mundo dos usuários da língua, inteligibilidade, atendimento ao solicitado no enunciado do ponto de vista do desenvolvimento do conteúdo, entre outros.
A grade de avaliação do desempenho linguístico considerou, portanto, três grandes grupos de competências, segundo os aspectos explicitados anteriormente : 1.
Aspectos ortográficos : domínio das convenções ortográficas : grafia de vogais e consoantes, uso de maiúsculas e minúsculas, emprego do hífen e acentuação gráfica ; 2.
Aspectos textuais : domínio dos procedimentos de estruturação textual do ponto de vista microestrutural : organização interna dos períodos, emprego de conectores para a articulação lógica entre os períodos e entre os parágrafos, emprego de marcas de referenciação lexical e pronominal ; utilização dos sinais de pontuação que contribuem para a organização lógica da frase. 3.
Aspectos morfossintáticos e vocabulares : domínio das regras de caráter morfossintático estabelecidas como modelares do ponto de vista da modalidade escrita formal da norma-padrão da Língua Portuguesa : concordância nominal e verbal, regência nominal e verbal, colocação pronominal, flexão nominal e verbal, correlação entre tempos e modos verbais, ausência de marcas de oralidade.
A seleção vocabular adequada à modalidade escrita formal da Língua Portuguesa foi incorporada a essa última competência, tendo em vista a intersecção entre as duas do ponto de vista das exigências do registro formal da modalidade escrita da norma-padrão.
Os resultados da avaliação correspondem aos seguintes aspectos observados em cada competência : Aspectos ortográficos : A correção foi realizada classificando os textos em cinco níveis, nível zero a quatro.
O desempenho dos estudantes ficou concentrado, majoritariamente, no nível 3 ( de 4 a 7 desvios ), devido, principalmente, aos desvios de acentuação.
O número de desvios de grafia foi reduzido.
Diferentemente do resultado do ENADE/2013, aumentou a porcentagem de textos sem desvios ( enquadrados no nível 4 ) e diminuiu a porcentagem de textos enquadrados no nível 2.
Alguns textos, com pior desempenho, foram enquadrados no nível 1 ( de 8 a 12 desvios ), enquanto o nível zero foi atribuído a pouquíssimos casos, já que ele revela falta absoluta de domínio das convenções ortográficas.
Observou -se, portanto, que existe uma diferença muito grande de desempenho nos dois aspectos analisados : baixo índice de desvios ortográficos e grande índice de desvios de acentuação.
Em vários casos, ocorre ausência completa de acentuação gráfica.
Os resultados revelam, portanto, que a tendência dominante entre os universitários brasileiros é a eliminação da acentuação gráfica, provavelmente motivada pela vivência dos jovens relacionada aos aplicativos de comunicação via internet ( redes sociais e e-mails ).
Nesse tipo de comunicação, devido ao ritmo intenso de troca de mensagens, o uso de acentos gráficos foi praticamente abolido.
Outro fator que pode ter relação com essa tendência é a ausência de esclarecimento dos meios de comunicação, das autoridades e das escolas sobre as decisões do Acordo Ortográfico da Língua Portuguesa de 1990, gerando um estado de indefinição para os estudantes.
Os casos mais sistemáticos de eliminação do acento indicador da sílaba tônica são :  palavras proparoxítonas ( “ ridiculos ”, “ publicas ”, “ lideres ”, “ politicos ” ) ;  palavras paroxítonas terminadas em ditongo crescente ( “ homicidios ”, “ latrocinio ”, “ individuo ”, “ dependencia ” ) ;  palavras oxítonas ( “ ninguem ”, “ esta ”, “ ate ”, “ ai ” ).
Por outro lado, destaca -se o uso indevido do acento gráfico em determinadas palavras, como observado nas grafias * “ jornáis ”, * “ telejornaís ”, * “ propíciar ”, * “ medídas ”, * “ dígnidade ”, * “ cídades ”.
Quanto ao domínio das convenções relativas à grafia das palavras, observam -se desvios como : a hipercorreção pela escolha de “ e ” no lugar de “ i ”, por influência de hábitos da oralidade ( * “ enumeros ” por “ inúmeros ”, * “ entevenção ” por “ intervenção ” ) ; a eliminação do “ r ” marcador do infinitivo verbal ( * “ esta ” no lugar de “ estar ” ).
Outros casos de desvios de grafia relacionados à variação diastrática podem ser observados em *“estrupo ”, *“automovís ”, *“viensse ”, *“camihada ”.
São muito frequentes os seguintes desvios de caráter ortográfico, com repercussão morfossintática :  eliminação da marca de infinitivo ( -r- ) e substituição por acento agudo ( “ está ” no lugar de “ estar ” ) ;  confusão entre “ ão ” e “ am ” nas formas verbais ( “ invadão ” no lugar de “ invadam ” e “ estam ” no lugar de “ estão ” no presente do indicativo ; “ estaram ” no lugar de “ estarão ” no futuro do indicativo ) ;  confusão entre a grafia do verbo “ haver ” ( “ há ” ) e o artigo definido ou a preposição “ a ” ;  uso de hífen para separar pronome átono – tanto uso indevido quanto omissão ( no pretérito imperfeito e futuro do subjuntivo : “ evitar -mos ” no lugar de “ evitarmos ”, “ percebesse ” no lugar de “ percebe -se ” e vice-versa ) ; Observam -se, também, muitos casos de inadequação no uso da maiúscula : ausência de diferença entre a primeira letra e as outras, em início de período, principalmente dos estudantes que adotam a escrita em letra de imprensa ; utilização de maiúscula para destacar determinadas palavras-chave do texto, como “ Violência ”, “ Brasileiros ”, “ Fatores ”, “ Ozônio ”, “ Sustentável ”.
Vale observar, também, que, ao contrário do que se esperava, não apareceram abreviaturas próprias do “ internetês ”, ou seja, dos hábitos de comunicação escrita adquiridos pelo uso de redes sociais e emails.
Aspectos textuais : Quanto a esses aspectos a correção também classificou os textos em cinco níveis ( zero a quatro ), em função da quantidade de erros apresentados.
O desempenho dos estudantes ficou concentrado, majoritariamente, nos níveis 3 e 2, devido à grande ocorrência de problemas de estruturação textual.
Foi muito baixo o número de textos que não apresentaram qualquer problema estrutural e, portanto, ficaram enquadrados no nível 4.
Alguns textos com pior desempenho foram enquadrados no nível 1, enquanto o nível zero foi atribuído a pouquíssimos casos, já que ele revela a existência de um texto sem articulação e com comprometimento do sentido.
Observou -se que a grande maioria dos estudantes não distribuiu as ideias em parágrafos, talvez devido ao pequeno número de linhas disponibilizadas para a resposta da questão ou, quem sabe, pela suposição de que não seria necessária essa divisão por não se tratar de um texto no modelo de uma redação dissertativo-argumentativa, como solicitado nos vestibulares.
Outro aspecto observado na estruturação textual foi a divisão em dois itens, provavelmente motivados pelo encaminhamento do enunciado das questões, que apresentavam dois pontos a serem detalhados : a questão 1 solicitava que os estudantes dissertassem sobre as consequências do transporte motorizado ( a ) e apresentassem ações de intervenção por parte do poder público ( b ) ; a questão 2 solicitava que os estudantes dissertassem sobre as causas da violência urbana ( a ) e os fatores para evitá -la ( b ).
Essa competência pode ser considerada como a mais problemática no que diz respeito ao desempenho linguístico dos estudantes, porque são muitos os problemas observados, desvios acumulados durante toda a formação escolar e que não se resolvem com um estudo autodidata, como acontece com regras ortográficas ou morfossintáticas : sequência justaposta de ideias sem encaixamentos sintáticos ; redução drástica de estruturas subordinadas, ao lado do aumento na frequência de estruturas coordenadas e absolutas ; redução no uso de conectores para expressar relações lógicas essenciais à construção do texto, substituídas pela exigência de inferência por parte do interlocutor para suprir a sua ausência ; emprego equivocado de operadores que não estabelecem relações lógicas coerentes entre ideias do texto ; emprego inadequado do pronome relativo ( com omissão da preposição ou a utilização de pronome inadequado, como “ onde ” ) ; repetição exaustiva de termos sem a utilização de procedimentos mais sofisticados de substituição ( hiperonímias, hiponímias, nominalizações, expressões metafóricas ) ; frases fragmentadas que comprometem a estrutura lógico-gramatical ; frases formadas apenas por oração subordinada, sem oração principal.
Um importante aspecto a destacar é o baixíssimo desempenho de uma parte dos estudantes em relação à estrutura formal do texto produzido, o que é extremamente preocupante ao se levar em conta que são graduandos em fase final de formação.
São frequentes os casos de desvios de estruturação frasal, com uso inadequado ou ausência de conectivos entre parágrafos e entre frases.
Em uma parte dos textos, falta um mínimo de textualidade e de domínio do registro padrão da língua.
Na verdade, observam -se relações linguísticas quase agramaticais, como as estabelecidas pela sequência de gerúndios sem o apoio de um ponto de partida para a organização das informações gramaticais e semânticas, ou seja, sem uma oração principal.
Quanto à utilização dos mecanismos de referenciação, deve -se destacar a ocorrência, em uma boa parte dos textos, de repetições de palavras ou expressões sem a utilização de termos sinônimos ou pronomes, como seria adequado.
Quanto à utilização dos sinais de pontuação, observou -se uma grande precariedade nos textos analisados.
É muito frequente a ocorrência de parágrafos sem marca interna de pontuação para separar os períodos.
Vale observar que não foi penalizada a ausência de vírgula para destacar locuções ou adjuntos adverbiais de pequena extensão deslocados de posição na frase, por ser um uso opcional.
São os seguintes os tipos de problemas encontrados : • vírgula : utilização de vírgula para separar o sujeito e o predicado ; ocorrência de apenas uma das vírgulas para separar uma palavra, uma expressão ou uma oração encaixada ; uso de vírgula no lugar do ponto para separar ideias que constituem períodos distintos ; ausência de vírgula para separar elementos de uma enumeração ; ausência de vírgula para separar oração adjetiva explicativa ou utilização inadequada para separar oração adjetiva restritiva ;  ponto e vírgula : utilização do ponto e vírgula no lugar de vírgula ;  ponto final : ausência de ponto final para separar períodos.
Aspectos morfossintáticos e vocabulares : Da mesma forma que nos aspectos anteriores os textos foram classificados em níveis de zero a quatro.
O desempenho dos estudantes ficou concentrado, majoritariamente, nos níveis 3 e 2, devido à grande ocorrência de problemas de regência e concordância.
O nível 4 foi atribuído a um número menor de textos.
Alguns textos com pior desempenho foram enquadrados no nível 1, enquanto o nível zero foi atribuído a pouquíssimos casos, já que ele revela a existência de um texto sem o respeito às mínimas exigências morfossintáticas da norma-padrão e com comprometimento do sentido.
Os resultados são muito transparentes em relação aos aspectos mais problemáticos do desempenho dos estudantes.
O desvio mais frequente, em relação à regência, é a falta do sinal indicativo da crase – isso revela que o usuário não tem consciência de que, sob a forma do termo “ a ”, existe a presença de uma preposição “ a ”, exigida pela regência do termo anterior.
Embora em outros exames, como o Enem, a falta de crase seja penalizada em acentuação, nesta correção esse desvio foi considerado no âmbito dos aspectos morfossintáticos.
Outro problema relacionado à regência verbal e à nominal, encontrado frequentemente nas questões, foi a ausência de preposição antes de pronome relativo, processo generalizado na modalidade oral da língua, em situações de registro informal.
Apesar da possibilidade de que essa alteração de regência se generalize no padrão escrito da Língua Portuguesa, como já está ocorrendo até em textos jornalísticos, o não emprego da preposição foi penalizado neste processo de avaliação.
Outro desvio muito frequente diz respeito aos processos de concordância verbal e de concordância nominal.
Quanto à concordância de número, observou -se ausência de marca ( com sujeito anteposto ou posposto ) ou uso indevido ( uso inadequado da marca de plural comandado pelo núcleo plural da locução adjetiva, apesar de o substantivo que funciona como núcleo do sintagma nominal estar no singular ).
Uma ocorrência que se destacou foi a ausência de acento circunflexo na forma plural do presente do indicativo dos verbos “ ter ” e “ vir ”, que foi considerada como um desvio na concordância verbal e não na acentuação gráfica.
Quanto à concordância de gênero, vários casos foram observados, normalmente no âmbito de sintagmas nominais longos, em que o adjetivo está afastado do substantivo.
Deve -se destacar uma ocorrência não observada no Enade/2013 : o aparecimento da marca de plural em verbos ou adjetivos relacionados a núcleos substantivos no singular, evidenciando um processo de hipercorreção.
Quanto à questão da colocação pronominal, foram poucos os casos observados.
Apesar de serem aspectos relacionados à oralidade, concluiu -se que, no registro escrito formal, a maioria dos estudantes já incorporou regras como a não introdução da frase por um pronome oblíquo e a próclise na presença de um termo atrator.
Não se adotou, entretanto, o padrão excessivamente formal descrito pelas gramáticas normativas em relação à posição do pronome oblíquo em locuções verbais, já que esse uso está muito distante da prática cotidiana, até em textos mais formais.
Quanto aos aspectos vocabulares, alguns tipos de inadequação foram observados : expressões da oralidade apareceram em algumas respostas, mas sem maior relevância do ponto de vista quantitativo ; seleção vocabular incompatível com o contexto, gerando situações de falta de inteligibilidade ; falta de domínio de vocabulário mais abstrato e de maior complexidade, essencial ao desenvolvimento do texto de base dissertativa.
Várias marcas de oralidade foram identificadas, embora não em alta frequência : o uso do pronome relativo “ onde ” como relativo universal, falta de artigo definido antes de substantivo, repetição de palavras por falta de vocabulário, reduções como “ tá ”, “ pra ”, “ pro ”, “ prum ”, expressões informais. 3.3.2 Componente de Conhecimento Específico Na parte da prova relativa às questões discursivas do Componente de Conhecimento Específico ( Tabela 3.15 ), observa -se que a média foi bem mais baixa do que para as questões discursivas do Componente de Formação Geral.
Enquanto no Componente de Formação Geral a média para estudantes de Pedagogia ( Licenciatura ) de todo o Brasil foi 50,7, na parte de Conhecimento Específico a média foi 24,6.
A maior média deste componente foi obtida pelos estudantes da região Sudeste ( 25,8 ), e a menor, pelos da região Centro-Oeste ( 22,6 ).
Quanto à variabilidade das notas, o desvio padrão de todo o Brasil foi 15,1.
O maior desvio concordância de gênero, vários casos foram observados, normalmente no âmbito de sintagmas nominais longos, em que o adjetivo está afastado do substantivo.
Deve -se destacar uma ocorrência não observada no Enade/2013 : o aparecimento da marca de plural em verbos ou adjetivos relacionados a núcleos substantivos no singular, evidenciando um processo de hipercorreção.
Quanto à questão da colocação pronominal, foram poucos os casos observados.
Apesar de serem aspectos relacionados à oralidade, concluiu -se que, no registro escrito formal, a maioria dos estudantes já incorporou regras como a não introdução da frase por um pronome oblíquo e a próclise na presença de um termo atrator.
Não se adotou, entretanto, o padrão excessivamente formal descrito pelas gramáticas normativas em relação à posição do pronome oblíquo em locuções verbais, já que esse uso está muito distante da prática cotidiana, até em textos mais formais.
Quanto aos aspectos vocabulares, alguns tipos de inadequação foram observados : expressões da oralidade apareceram em algumas respostas, mas sem maior relevância do ponto de vista quantitativo ; seleção vocabular incompatível com o contexto, gerando situações de falta de inteligibilidade ; falta de domínio de vocabulário mais abstrato e de maior complexidade, essencial ao desenvolvimento do texto de base dissertativa.
Várias marcas de oralidade foram identificadas, embora não em alta frequência : o uso do pronome relativo “ onde ” como relativo universal, falta de artigo definido antes de substantivo, repetição de palavras por falta de vocabulário, reduções como “ tá ”, “ pra ”, “ pro ”, “ prum ”, expressões informais. 3.3.2 Componente de Conhecimento Específico Na parte da prova relativa às questões discursivas do Componente de Conhecimento Específico ( Tabela 3.15 ), observa -se que a média foi bem mais baixa do que para as questões discursivas do Componente de Formação Geral.
Enquanto no Componente de Formação Geral a média para estudantes de Pedagogia ( Licenciatura ) de todo o Brasil foi 50,7, na parte de Conhecimento Específico a média foi 24,6.
A maior média deste componente foi obtida pelos estudantes da região Sudeste ( 25,8 ), e a menor, pelos da região Centro-Oeste ( 22,6 ).
Quanto à variabilidade das notas, o desvio padrão de todo o Brasil foi 15,1.
O maior desvio padrão foi encontrado nas regiões Norte e Nordeste ( 15,5 ), e o menor, na região Sul ( 14,5 ).
A maior nota máxima foi obtida na região Sul ( 90,0 ), enquanto a menor nota máxima foi encontrada na região Norte ( 76,7 ).
Além disso, a nota mínima ( 0,0 ) foi obtida por pelo menos um estudante em todas as regiões do Brasil, sem exceção.
A mediana do Brasil como um todo foi 25,0, a mesma obtida na região Sul.
Nas regiões Norte e Nordeste a mediana foi 23,3, e nas regiões Sudeste e Centro-Oeste foi 26,7 e 21,7, respectivamente.
O Gráfico 3.10 representa a distribuição das notas nas questões discursivas no Componente de Conhecimento Específico.
A moda desta distribuição ocorre no intervalo ( 20 ; 30 ].
Nota -se que a distribuição após a moda é decrescente a partir do intervalo modal.
A análise de cada uma destas questões será feita a seguir. 3.3.2.1 Análise da Questão Discursiva 3 do Componente de Conhecimento Específico Na questão 3, cujos resultados aferidos encontram -se descritos na Tabela 3.16, a média dos estudantes de todo o Brasil foi 27,9.
A menor média nessa questão foi obtida pelos estudantes da região Norte ( 25,3 ), enquanto a maior média foi obtida na região Sudeste ( 29,1 ).
Quanto à variabilidade das notas, o desvio padrão de todo o Brasil foi 20,4.
O maior desvio padrão foi obtido na região Sudeste ( 20,7 ), enquanto o menor foi obtido na região Centro-Oeste ( 19,8 ).
A nota máxima, 100,0 pontos, foi alcançada por, pelo menos, um estudante em todas as regiões do Brasil, sem, exceção, assim como a nota mínima ( 0,0 ).
A mediana do Brasil como um todo foi 30,0, a mesma obtida nas regiões Sudeste e Sul.
Nas demais regiões a mediana foi 25,0.
O Gráfico 3.11 mostra a distribuição das notas na questão discursiva 3, do Componente de Conhecimento Específico da área de Pedagogia ( Licenciatura ).
Essa distribuição apresenta a moda principal no intervalo ( 20 ; 30 ]. 3.3.2.2 Comentários sobre as respostas à Questão Discursiva 3 O enunciado da questão trazia dois pequenos fragmentos de textos.
O primeiro era a citação de uma frase de Larrosa ( 2004 ) sobre infância, e o segundo, um relato de situação de aula protagonizada por uma professora.
A frase que introduzia o comando da questão era iniciada por trecho em que se destacavam os modos peculiares de percepção do mundo pelas crianças, seguido de considerações acerca do papel da sala de aula como espaço privilegiado para a “ troca e ampliação de saberes e conhecimentos ”.
O comando em si solicitava dos estudantes a apresentação de uma proposta de trabalho a ser desenvolvida com a turma regida pela professora que provocara a situação mencionada.
Essa proposta deveria abordar os aspectos solicitados em três itens.
O item ‘ a ’, solicitava que os estudantes elencassem elementos constitutivos do planejamento “ que levem em conta as experiências das crianças ”.
Parte expressiva dos examinados respondeu a esse item destacando a necessidade do respeito às experiências trazidas pelas crianças, ao conhecimento prévio sobre os temas objeto da aula e à interpretação que fazem da realidade, mesmo que imaginária.
Contudo, poucos mencionaram a importância de uma avaliação diagnóstica ou usaram este termo associado às respectivas estratégias de elaboração do ato de planejar a ação e prever os resultados esperados.
Outro dado relevante foi o fato de que poucas respostas indicaram a necessidade de \textbf{traçar} objetivos para alcançar os resultados desejados, o que poderia revelar a inobservância dos riscos decorrentes da desconexão entre o planejamento e a execução da ação pedagógica.
No tocante ao item ‘ b ’, que solicitava a proposição de atividades a serem realizadas, considerando -se uma perspectiva interdisciplinar, pôde -se observar maior quantidade de respostas adequadas ou próximas ao que foi estabelecido como padrão.
A maioria dos examinados descreveu, com riqueza de detalhes, uma variedade de atividades que poderiam ser desenvolvidas pela professora, no sentido de trabalhar o conteúdo e as competências relativas ao focalizado pelo enunciado.
Entre as atividades propostas, registraram -se, por exemplo, visitas orientadas ao campo e à cidade, construção de maquetes, elaboração de mapas e gráficos, rodinhas de conversa em sala de aula, dramatizações, passeios ou excursões a fábricas, sítios ou fazendas, entrevistas, questionários, pesquisas em geral.
Por fim, o item ‘ c ’ solicitava a indicação de “ formas de registro para avaliação das crianças ”.
As respostas dadas a esse item foram as que apresentaram maior inconsistência, pois evidenciaram uma quase completa ausência de preocupação do estudante de Pedagogia em conciliar instrumentos específicos de avaliação às atividades propostas, o que, por certo, se relaciona à baixa associação entre a realização do planejamento ( inclusive, com a elaboração de objetivos definidos ) e o desenvolvimento e execução das atividades.
Importa, contudo, ressaltar que os examinados, em grande parte, reconheceram a necessidade de avaliação do processo de aprendizagem, de preferência quando implementada de modo permanente e contínuo.
Técnicas como a de observação da fala e da produção das crianças foram sugeridas, mas raramente se especificou, de forma clara, os instrumentos a serem empregados para esse fim.
Os poucos estudantes que assim o fizeram, sugeriram o uso de relatórios, diários de bordo, registros de análise de desenhos e portfólios.
O padrão de respostas apresentou três subdivisões, cada qual voltada para um dos aspectos específicos em que a questão se desdobrava.
Os pontos centrais a serem observados na correção da questão foram transmitidos de forma clara, com linguagem precisa.
O padrão assegurava um grau adequado de flexibilidade que permitiu considerar a diversidade de reflexões encontradas nas respostas, com diferentes níveis de profundidade e de especificidade em relação ao conteúdo da questão.
Pode -se considerar que a questão se situa no mesmo nível de dificuldade conceitual das demais questões discursivas de componente específico, embora, apresente algumas particularidades relacionadas ao fato de exigir uma resposta mais “ prática ”, aplicada ao planejamento, condução e avaliação do processo pedagógico.
Apesar de o padrão de respostas contemplar um amplo conjunto de possibilidades de solução, os estudantes dos cursos de Pedagogia, de uma forma geral, não apresentaram bom rendimento, em especial nos itens ‘ a ’ e ‘ c ’.
Tal constatação pode estar relacionada à incompreensão do enunciado, a dificuldades de redação de texto argumentativo, como prescrevia o enunciado, ou, ainda, a lacunas na formação para o magistério.
Foram muito frequentemente registrados erros de encadeamento lógico nos textos produzidos.
Evidenciou -se, também, o precário domínio dos aspectos teórico-metodológicos envolvidos nas fases do planejamento da ação pedagógica.
Foi realizada uma classificação das respostas em três níveis, em função das notas obtidas : fracas ( notas de zero a 30 ) ; medianas ( notas de 35 a 70 ) e boas ( notas de 75 a 100 ).
A maior parte das notas se concentrou nas faixas consideradas fracas ou medianas, com predominância de notas entre 30 e 40 pontos.
Houve poucos O estudante de Pedagogia em conciliar instrumentos específicos de avaliação às atividades propostas, o que, por certo, se relaciona à baixa associação entre a realização do planejamento ( inclusive, com a elaboração de objetivos definidos ) e o desenvolvimento e execução das atividades.
Importa, contudo, ressaltar que os examinados, em grande parte, reconheceram a necessidade de avaliação do processo de aprendizagem, de preferência quando implementada de modo permanente e contínuo.
Técnicas como a de observação da fala e da produção das crianças foram sugeridas, mas raramente se especificou, de forma clara, os instrumentos a serem empregados para esse fim.
Os poucos estudantes que assim o fizeram, sugeriram o uso de relatórios, diários de bordo, registros de análise de desenhos e portfólios.
O padrão de respostas apresentou três subdivisões, cada qual voltada para um dos aspectos específicos em que a questão se desdobrava.
Os pontos centrais a serem observados na correção da questão foram transmitidos de forma clara, com linguagem precisa.
O padrão assegurava um grau adequado de flexibilidade que permitiu considerar a diversidade de reflexões encontradas nas respostas, com diferentes níveis de profundidade e de especificidade em relação ao conteúdo da questão.
Pode -se considerar que a questão se situa no mesmo nível de dificuldade conceitual das demais questões discursivas de componente específico, embora, apresente algumas particularidades relacionadas ao fato de exigir uma resposta mais “ prática ”, aplicada ao planejamento, condução e avaliação do processo pedagógico.
Apesar de o padrão de respostas contemplar um amplo conjunto de possibilidades de solução, os estudantes dos cursos de Pedagogia, de uma forma geral, não apresentaram bom rendimento, em especial nos itens ‘ a ’ e ‘ c ’.
Tal constatação pode estar relacionada à incompreensão do enunciado, a dificuldades de redação de texto argumentativo, como prescrevia o enunciado, ou, ainda, a lacunas na formação para o magistério.
Foram muito frequentemente registrados erros de encadeamento lógico nos textos produzidos.
Evidenciou -se, também, o precário domínio dos aspectos teórico-metodológicos envolvidos nas fases do planejamento da ação pedagógica.
Foi realizada uma classificação das respostas em três níveis, em função das notas obtidas : fracas ( notas de zero a 30 ) ; medianas ( notas de 35 a 70 ) e boas ( notas de 75 a 100 ).
A maior parte das notas se concentrou nas faixas consideradas fracas ou medianas, com predominância de notas entre 30 e 40 pontos.
Houve poucos registros de respostas classificadas como boas : poucas respostas dadas a essa questão mereceram notas superiores a 70 pontos.
Entre os examinados, um pouco mais de 10 \% deixou essa questão em branco, e quase 5 \% deram respostas às quais foram atribuídas nota zero.
Conforme já sinalizado, contatou -se que, em geral, as respostas centraram -se na apresentação e descrição de atividades, evidenciando uma provável desconexão entre os atos de planejar, de executar e de avaliar a ação pedagógica.
Dessa forma, parece haver a predominância de certo “ ativismo pedagógico ”, identificado no expressivo número de respostas cujo foco se voltou exclusivamente para a proposição de atividades de desenvolvimento do conteúdo, sem, contudo, estabelecer a indispensável vinculação com o planejamento e a avaliação do conteúdo proposto.
Os resultados obtidos na presente análise demonstram que a formação destes profissionais demanda maiores cuidados.
Faz -se necessário que os cursos promovam uma maior reflexão sobre a função do pedagogo no papel de estrategista e articulador da ação pedagógica, junto aos demais atores institucionais.
É preocupante constatar que prevalece entre os estudantes dos cursos de Pedagogia a noção de que basta propor atividades de desenvolvimento do conteúdo, sem que sejam necessários o planejamento, precedido de diagnóstico, e a avaliação, para que o processo de ensino-aprendizagem ocorra de forma satisfatória. 3.3.2.3 Análise da Questão Discursiva 4 do Componente de Conhecimento Específico A Tabela 3.17 contém as informações relativas à questão 4 do conjunto de questões do Componente de Conhecimento Específico.
O desempenho dos estudantes nessa questão foi inferior ao das questões 3 e 5.
A média geral do Brasil foi 18,5, sendo a menor média registrada na região Norte ( 16,7 ), e a maior, na região Sudeste ( 19,5 ).
A nota máxima ( 100,0 ) foi atingida por, pelo menos, um concluinte em todas as regiões do Brasil, exceto na região Norte ( 95,0 ).
A mediana em todo o Brasil foi 15,0, o mesmo valor foi obtido nas regiões Sudeste e Sul, enquanto nas demais regiões a mediana foi 10,0.
As notas mínimas foram zero em todas as regiões, sem exceção. registros de respostas classificadas como boas : poucas respostas dadas a essa questão mereceram notas superiores a 70 pontos.
Entre os examinados, um pouco mais de 10 \% deixou essa questão em branco, e quase 5 \% deram respostas às quais foram atribuídas nota zero.
Conforme já sinalizado, contatou -se que, em geral, as respostas centraram -se na apresentação e descrição de atividades, evidenciando uma provável desconexão entre os atos de planejar, de executar e de avaliar a ação pedagógica.
Dessa forma, parece haver a predominância de certo “ ativismo pedagógico ”, identificado no expressivo número de respostas cujo foco se voltou exclusivamente para a proposição de atividades de desenvolvimento do conteúdo, sem, contudo, estabelecer a indispensável vinculação com o planejamento e a avaliação do conteúdo proposto.
Os resultados obtidos na presente análise demonstram que a formação destes profissionais demanda maiores cuidados.
Faz -se necessário que os cursos promovam uma maior reflexão sobre a função do pedagogo no papel de estrategista e articulador da ação pedagógica, junto aos demais atores institucionais.
É preocupante constatar que prevalece entre os estudantes dos cursos de Pedagogia a noção de que basta propor atividades de desenvolvimento do conteúdo, sem que sejam necessários o planejamento, precedido de diagnóstico, e a avaliação, para que o processo de ensino-aprendizagem ocorra de forma satisfatória. 3.3.2.3 Análise da Questão Discursiva 4 do Componente de Conhecimento Específico A Tabela 3.17 contém as informações relativas à questão 4 do conjunto de questões do Componente de Conhecimento Específico.
O desempenho dos estudantes nessa questão foi inferior ao das questões 3 e 5.
A média geral do Brasil foi 18,5, sendo a menor média registrada na região Norte ( 16,7 ), e a maior, na região Sudeste ( 19,5 ).
A nota máxima ( 100,0 ) foi atingida por, pelo menos, um concluinte em todas as regiões do Brasil, exceto na região Norte ( 95,0 ).
A mediana em todo o Brasil foi 15,0, o mesmo valor foi obtido nas regiões Sudeste e Sul, enquanto nas demais regiões a mediana foi 10,0.
As notas mínimas foram zero em todas as regiões, sem exceção.
O Gráfico 3.12 representa a distribuição de notas da questão discursiva 4, do Componente de Conhecimento Específico.
Essa distribuição apresenta a moda principal nas notas zero, seguida de perto pelos intervalos ( 10 ; 20 ] e ( 20 ; 30 ]. 3.3.2.4 Comentários sobre as respostas à Questão Discursiva 4 A questão solicitava ao estudante redigir um texto dissertativo a respeito do tema “ a escola cidadã e a construção da cidadania ”, problemática contemporânea da formação dos profissionais da área de Pedagogia, conforme verificável pela leitura das Diretrizes Curriculares da Área.
O enunciado trazia dois textos com definições que deveriam ser consideradas na redação do estudante.
Um desses textos definia “ comunidade escolar ”, e o outro, “ movimentos sociais ”.
Além disso, o comando da questão orientava a redação do texto dissertativo por meio de três itens, os quais indicavam conceitos e relações que deveriam ser considerados : conceito de escola cidadã ( item ‘ a ’ ) ; relação da escola cidadã com a construção da cidadania ( item ‘ b ’ ) ; e diálogo estabelecido entre escola, comunidade e movimentos sociais ( item ‘ c ’ ).
O enunciado da questão era claro, embora os textos motivadores contribuíssem de forma mais direta e imediata somente para o atendimento ao solicitado pelo item ‘ c ’ da questão.
Comparando a questão com as demais questões discursivas de Conhecimento Específico, podemos considerá -la como uma questão “ difícil ”, por exigir desenvoltura no campo conceitual, enfatizada, sobretudo, em disciplinas como “ sociologia da educação ”, “ políticas educacionais ” e “ filosofia da educação ”.
O padrão de respostas utilizado para a correção da questão atendia aos critérios de pertinência e exaustividade, dando conta da variedade de enfoques em relação a cada um dos itens em que a questão se dividia.
Para cada item o padrão indicava, por meio de exemplos, um conjunto de alternativas de respostas com nível de profundidade mediano, o que assegurou a necessária flexibilidade da correção.
Percebeu -se que as respostas poderiam ser classificadas quanto aos diferentes tipos de conteúdo, independente de estes estarem ou não corretos ( em sua grande maioria não estavam ).
É mesmo possível supor, dada a similaridade das respostas, que tenha existido algum tipo de preparação para o exame, análoga à praticada pelos cursos preparatórios para certames públicos.
Quanto à abordagem do tema adotada pelos estudantes, pôde -se notar que, com frequência, eram reproduzidas percepções eivadas de preconceitos calcados no senso comum.
Tal constatação indica que a formação universitária para a maioria não enfatizou suficientemente, o estímulo ao pensamento crítico, baseado na pesquisa e na reflexão.
A correção das respostas revelou um quadro de incompreensões, do qual se destaca uma percepção “ vaga ” do conceito de cidadania.
Esta é, muitas vezes, confundida com melhoria de vida ; e o exercício da cidadania, considerado exclusivamente como o direito de dispor de serviços públicos básicos.
Na verdade, para muitos dos estudantes a expressão “ cidadão ” é compreendida como sinônimo de “ indivíduo ”, um desvio bastante presente no senso comum brasileiro, que remonta, talvez, ao início do período republicano.
Isso demonstra que, apesar da permanência em curso superior da área das Ciências Humanas, esta diferenciação, elementar e necessária, não foi devidamente compreendida pelo estudante.
O “ cidadão ” é tomado como o indivíduo honesto, portanto a formação para a cidadania é confundida com a incorporação de elementos morais, em particular a honestidade.
Daí o uso amplamente empregado de expressões como “ cidadão de bem ”, “ cidadão de caráter ” ou até mesmo “ cidadania de caráter ”.
Assim, tal deformação conceitual, reduz a noção de cidadania a uma esfera estritamente comportamental, no sentido de um modelo de conduta.
Ainda nesse sentido, foi significativo o fato de, em muitas respostas, o estudante ter sublinhado a palavra “ respeito ”, em textos que sinalizam a presença de situações conflituosas, não raro violentas, dentro da escola.
Em função disso é possível entender porque, para muitos estudantes, “ formar cidadãos ” é entendido como “ formar indivíduos que respeitam os outros ”.
Deparamo -nos, em segundo lugar, com grande dificuldade de exposição do pensamento e explicitação de argumentos.
Talvez como resultado do acesso apenas superficial ao debate contemporâneo travado no campo da Educação, percebe -se que houve apenas memorização acrítica de expressões, mal compreendidas.
Assim, elas terminam sendo empregadas como “ chavões ” ou “ frases feitas ”, sem denotar maior domínio acerca das categorias empregadas no texto produzido como resposta.
Predominaram as respostas que denotavam não somente conhecimento superficial da temática, mas, também, uma “ antropomorfização ” de conceitos e categorias analíticas.
Por exemplo, conceitos como “ cidadania ” ou “ movimentos sociais ” foram, muitas vezes, tratados como se fossem entes vivos e não construções abstratas, voltadas para a compreensão da realidade social em geral e, particularmente, do universo escolar.
Como consequência, nos deparamos com o uso de diversas expressões, absorvidas acriticamente das Ciências Sociais e da Teoria Política, usadas sob a forma de meros jargões.
São exemplos : “ educação de qualidade ”, “ estudante integral ”, “ cidadão crítico e reflexivo ”, entre outras.
De maneira semelhante, a expressão “ escola cidadã ”, foi frequentemente considerada como “ utopia ”, um estágio ideal, porém inalcançável, geralmente por força do agir dos estudantes e de seus pais.
Quanto à participação da comunidade, o foco principal se voltou para os pais, e a participação foi quase sempre entendida como o acompanhamento das atividades escolares dos filhos.
A escassa participação dos pais foi apontada como um dos grandes problemas da Educação, fator responsável pelo insucesso escolar, não só, mas principalmente, nas redes públicas de ensino.
Caberia, ainda, registrar um equívoco muito frequentemente observado nas respostas : o de confundir “ opinar ” ou “ sugerir ” com “ decidir ”.
Nesse sentido, foi possível identificar demonstrações de rejeição à participação de outros segmentos da comunidade escolar que não os constituídos por funcionários da escola.
Finalmente, conviria assinalar uma preocupante tendência identificada nas respostas, expressa na ideia, compartilhada por boa parte dos estudantes, de que a escola e a sociedade estão separadas, constituindo dois espaços distintos e autônomos.
Persiste uma visão de que as crianças vivem inicialmente apenas nos espaços familiar e da escola, e que somente após deixar a “ vida escolar ” passariam a integrar a sociedade, o que consiste um equívoco sociológico primário.
O desempenho dos estudantes na questão foi ainda mais fraco do que o registrado em relação à questão 3, provavelmente, por exigir incursões mais profundas no plano conceitual.
Nessa questão cerca de 15 \% dos estudantes deixaram o espaço para a resposta “ em branco ”.
Foi perceptível que especificamente ao empregar a expressão “ compreensão da realidade social ” – muito assídua nas respostas – os examinados quase sempre o fizeram em um sentido normativo, de prescrição de direitos e deveres previamente estabelecidos.
Detectou -se falta de um pensamento minimamente analítico, realizado com base em conceitos científicos e categorias filosóficas.
Este caráter normativo, implícito em quase todas as respostas, implicava atribuir à escola a responsabilidade por assegurar o conhecimento de direitos e deveres pelos indivíduos, e, com efeito, por formar “ cidadãos de bem ”.
É preciso destacar, mais uma vez, que essa expressão – cidadãos de bem –, muito encontrada nos textos produzidos, raramente era acompanhada de quaisquer comentários a respeito do seu significado, sinalizando para a força que concepções oriundas do senso comum ocupam no discurso do estudante ; concepções que, ao que parece, dificilmente teriam sido objeto de maior problematização ao longo da formação superior dos examinados.
Assim sendo, o exercício da cidadania foi, não raramente, interpretado como o esforço de doutrinação por parte da escola, tanto em relação ao conjunto dos estudantes, quanto, mais amplamente, à “ comunidade escolar ”.
A escola, reproduzindo antigos ideários modernizadores dos anos quarenta a setenta, deveria agir sobre a comunidade para “ transformá -la ”.
Respostas em que se sublinhou a importância de relações mais “ horizontais ” no processo institucional de tomada de decisão foram, por consequência, bem mais raras.
Embora em sentido distinto, mas também igualmente de teor prescritivo, a cidadania, por muitas vezes, foi identificada como uma espécie de “ ajuda mútua ” entre os integrantes da comunidade escolar, visão muito relacionada ao entendimento da escola cidadã como uma “ escola aberta ”, ou seja, um local onde ocorreriam atividades distintas daqueles mais comuns ao processo de ensino e aprendizagem.
Muitas vezes, a identificação da “ escola cidadã ” como “ escola aberta ” se fazia acompanhada da presença de uma falsa antinomia entre “ ação escolar que enfatiza aprendizado ” e “ ação pedagógica que favorece o conhecimento de direitos e deveres ”, geralmente considerando como mais positiva a escola que privilegiasse o “ conhecimento de direitos e deveres ” em detrimento do aprendizado dos conteúdos das disciplinas escolares.
Seja pelas graves lacunas conceituais apresentadas, seja pela preocupante dificuldade de expressão exibida, conclui -se que os futuros profissionais da área da Educação têm recebido uma formação muito aquém da necessária para atuar no magistério, tanto da Educação Infantil quanto dos primeiros anos do Ensino Fundamental, momentos cruciais na escolarização de todo indivíduo. 3.3.2.5 Análise da Questão Discursiva 5 do Componente de Conhecimento Específico A Tabela 3.18 contém as informações relativas à questão 5 do conjunto do Componente de Conhecimento Específico.
A nota média dos estudantes de todo o Brasil foi 27,3.
A maior média foi registrada na região Sudeste ( 29,0 ), enquanto a menor média foi registrada na região Centro-Oeste ( 24,8 ).
Quanto à variabilidade das notas, o desvio padrão dos estudantes do Brasil, como um todo, foi 24,1.
Enquanto o maior desvio foi encontrado na região Norte ( 24,7 ), o menor foi encontrado na região Centro-Oeste ( 23,2 ).
A mediana do Brasil, como um todo, foi 30,0, a mesma obtida na região Sudeste, enquanto nas demais regiões a mediana foi 20,0.
A nota máxima ( 100,0 ) e a nota mínima ( 0,0 ) são as mesmas em todas as regiões, sem exceção.
O Gráfico 3.13 representa a distribuição das notas da questão discursiva 5 do Componente de Conhecimento Específico.
Destaca -se, o grande número de estudantes que deixaram esta questão em branco, correspondendo à moda da distribuição. 3.3.2.6 Comentários sobre as respostas à Questão Discursiva 5 O enunciado da questão era claro quanto à linguagem utilizada e não dava margem a interpretações diferentes da esperada.
O enunciado trazia a citação do trecho do texto legal da Resolução CNE/CP n{\EmojiFont °}01, de 15/05/2006, que caracteriza as habilitações dos formados em Licenciaturas em Pedagogia.
Infere -se que ao buscar articular esse trecho com o tema proposto para o texto dissertativo – que deveria enfatizar a atuação do Pedagogo – os estudantes tenham acabado por usar os termos ” professor ” e “ pedagogo ” de modo indistinto.
A questão solicitava que o estudante redigisse uma dissertação sobre “ A atuação do pedagogo em espaços não escolares ”, exigindo que fossem abordados aspectos desmembrados em dois itens.
O item ‘ a ’ instava o estudante a citar exemplos de atividades que poderiam ser desenvolvidas por um pedagogo num espaço não escolar escolhido pelo respondente.
Solicitava, ainda, que o texto contemplasse a relação entre essas atividades e aspectos teóricos da formação desse profissional.
Observa -se que o item destinava -se a orientar parte da redação da resposta, mas de forma bastante aberta, admitindo, portanto, diversas possibilidades de caminhos e abordagens.
O item ‘ b ’ orientava o concluinte a defender “ a importância da participação do pedagogo em espaços não escolares ”.
Esperava -se que essa defesa se relacionasse, pelo menos, com a atividade escolhida como resposta ao item ‘ a ’, ou seja, presumia -se ter como resposta a menção de pelo menos um espaço não escolar, seguida da descrição das atividades que poderiam ser realizadas nesse espaço, e a referência aos aspectos teóricos envolvidos.
Esperava -se, ainda, que fossem mencionadas funções do pedagogo em espaços não escolares e sua interação com demais profissionais que lá atuariam.
O mais marcante, numa análise geral das respostas dadas para a questão, foi a pouca consistência teórica e prática.
Além disso, da mesma forma que relatado em relação às outras duas questões discursivas, evidenciaram -se as dificuldades que os estudantes dos cursos de Pedagogia apresentam na exposição de suas ideias e no uso adequado da Língua Portuguesa.
Os erros ortográficos e gramaticais presentes na escrita, que, por si sós, já comprometem a compreensão, eram agravados por redações com pouca coesão e coerência de ideias.
Dessa forma, a exposição dos conhecimentos, mesmo que tenham sido construídos, se torna muito comprometida.
O padrão de respostas foi formulado de maneira abrangente, contemplando grande diversidade de respostas, o que possibilitou identificar nos exemplos de atividades escolhidos, sua coerência com o espaço não escolar mencionado e com aspectos teóricos, os quais, vale destacar foram pouco mencionados de forma adequada.
O critério de pontuação adotado subdividia o total de pontos de cada item, de modo a ser possível pontuar cada menção.
Poucos estudantes atenderam plenamente o esperado pelo padrão de respostas estabelecido, ou seja, poucos mencionaram espaços não escolares de atuação do pedagogo, descreveram atividades desenvolvidas nos espaços citados, relacionaram atividade com fundamentação teórica do curso e expressaram interação com outras áreas, defendendo de forma adequada a atuação do pedagogo.
A maior parte dos examinados apenas citava um espaço não escolar ( nem sempre correto ) e fazia menção a atividades que poderiam ser desenvolvidas nesse mesmo espaço.
Muitos apenas indicaram espaços não escolares, muitas vezes evidenciando a incapacidade de distinguirem espaços não escolares de espaços extraclasses.
Não Não poucos deram respostas incompatíveis com o que foi perguntado, e merece destaque a pobreza de vocabulário empregado.
O desempenho na questão foi parecido com o registrado em relação à questão 3.
No entanto, a quantidade de respostas deixadas em branco foi um pouco maior do que a verificada nas questões anteriores.
O percentual de respostas anuladas ou desconsideradas também foi elevado.
Observamos que um expressivo número de formandos em Pedagogia chegou ao final do curso não diferenciando o pedagogo e o professor, como também demonstraram não distinguir espaço escolar de espaço não escolar, de atuação do pedagogo.
Como já citado, um elevado número de examinados concebe espaços não escolares como se tratando de espaços extraclasses.
Pode -se deduzir, pelas respostas dos examinados, que foram poucos aqueles que chegaram a trabalhar e/ou estagiar em espaços não escolares de atuação do pedagogo.
Destaca -se, ainda, o fato de que o conhecimento de aspectos técnicos, didáticos e metodológicos, relacionados com a elaboração, desenvolvimento e avaliação de projetos, foram no máximo apenas tangenciados.
Muitos examinados demonstraram desconhecer o trabalho do pedagogo em espaços não escolares.
Foi perceptível que a causa da falta de respostas a muitos aspectos abarcados na referida questão deveu -se ao desconhecimento real sobre tema.
Tal fato pode estar relacionado aos currículos de muitos cursos de Pedagogia, que não incluem – ou incluem de maneira restrita – conteúdo disciplinar sobre espaços não escolares em seus programas.
A revelação mais imediata e que causa bastante preocupação é a baixa qualidade da formação dos estudantes dos cursos de Pedagogia.
Essa questão, torna ainda mais clara a falta de articulação da prática com a teoria, já evidenciada pelas respostas dadas às questões 3 e 4.
Traz à tona a constatação de que a diversidade de habilitações dos formandos de cursos de Licenciatura em Pedagogia, citada explicitamente no caput dessa questão, constitui nada mais que um objetivo, para o alcance do qual os cursos de Pedagogia não têm dado uma efetiva contribuição.
A questão 3, por força de sua formulação, deu ensejo a respostas em que determinadas práticas do processo de ensino-aprendizagem são percebidas, acriticamente, como “ naturais e universais ”, como se as possibilidades de acompanhamento e atenção por parte dos responsáveis pelos estudantes dos primeiros anos do ensino fundamental não viessem a ser sobredeterminadas por elementos de caráter social, como a propriedade de classe.
Também se confirmou, pela avaliação do conjunto total de respostas relativas à questão 5, o entendimento equivocado, demonstrado por grande parte dos examinados, de que a atuação profissional dos pedagogos pode ocorrer em quaisquer campos de trabalho, independente dos aspectos constituintes dos mesmos, considerando -se o pedagogo como um profissional adaptável a todos os processos e condições de trabalho.
As respostas dos examinados sinalizaram para percepções empiristas dos processos educacionais, com baixa compreensão analítica dos mesmos.
Assim, se o possível intento dos formuladores das questões era possibilitar respostas que manifestassem diversificada compreensão teórica, isto foi sumamente contrariado pelo teor predominante das respostas apresentadas.
Para finalizar, apesar de já bastante comentado, é preocupante que profissionais da Educação – futuros formadores polivalentes de estudantes de diversos níveis, futuros planejadores de ações, gestores, supervisores ou dirigentes educacionais, entre tantas atuações que podem vir a ser assumidas pelos estudantes que prestaram o ENADE/2014 – evidenciem tanta dificuldade para se comunicar na forma escrita.

% matched lemmas: ajudar, caminhada, carbono, dióxido, traçar
\end{textsample}
